\documentclass[11pt]{article}
\usepackage[english]{babel}
\usepackage{amsmath,amssymb}
\usepackage[colorlinks=true]{hyperref}
\hypersetup{urlcolor=blue, citecolor=red, linkcolor=blue}
\usepackage{amsthm}

\usepackage{yhmath} 

\usepackage[T1]{fontenc}

\usepackage[a4paper, left=25mm, right=25mm, top=30mm, bottom=30mm]{geometry}

\usepackage{mathtools}
\usepackage{mathrsfs}

%%%%%%%%%% Start TeXmacs macros
\newcommand{\assign}{:=}
\newcommand{\nobracket}{}
\newcommand{\tmem}[1]{{\em #1\/}}
\newcommand{\tmop}[1]{\ensuremath{\operatorname{#1}}}
\newcommand{\tmrsup}[1]{\textsuperscript{#1}}
\newcommand{\tmtextbf}[1]{\text{{\bfseries{#1}}}}
\newcommand{\tmtextit}[1]{\text{{\itshape{#1}}}}
\newcommand{\tmtextmd}[1]{\text{{\mdseries{#1}}}}
\newcommand{\tmtextrm}[1]{\text{{\rmfamily{#1}}}}
\newcommand{\tmtextsc}[1]{\text{{\scshape{#1}}}}
\newcommand{\tmtexttt}[1]{\text{{\ttfamily{#1}}}}
\newcommand{\tmtextup}[1]{\text{{\upshape{#1}}}}


\theoremstyle{plain}
\newtheorem{theorem}{Theorem}[section]
\newtheorem{corollary}[theorem]{Corollary}
\newtheorem{proposition}[theorem]{Proposition}
\newtheorem{lemma}[theorem]{Lemma}

\theoremstyle{definition}
\newtheorem{definition}[theorem]{Definition}
%\newtheorem{note}[theorem]{Note}
\newtheorem{remark}[theorem]{Remark}

\newcommand\xqed[1]{%
    \leavevmode\unskip\penalty9999 \hbox{}\nobreak\hfill
	\quad\hbox{#1}}
\newcommand\demo{\xqed{$\triangle$}}


\newtheorem{example}[theorem]{Example}


% aplicacio restringida
\newcommand\restr[2]{{% we make the whole thing an ordinary symbol
  \left.\kern-\nulldelimiterspace % automatically resize the bar with \right
  #1 % the function
  %\vphantom{\big|} % pretend it's a little taller at normal size
  \right|_{#2} % this is the delimiter
}}

% matriu transposada
\newcommand*{\transp}[2][-3mu]{\ensuremath{\mskip1mu\prescript{\smash{\mathrm t\mkern#1}}{}{\mathstrut#2}}}%

% longhookrightarrow
\newcommand*{\longhookrightarrow}{\ensuremath{\lhook\joinrel\relbar\joinrel\rightarrow}}
% \newcommand{\me}{\mathrm{e}}
% \newcommand{\sq}{\mathrm{sq}}
% \newcommand{\simp}{\mathrm{Simp}^{or}}
% \newcommand{\sing}{\mathrm{Sing}}
% \newcommand{\N}{\mathbb{N}}
% \newcommand{\Z}{\mathbb{Z}}
% \newcommand{\Q}{\mathbb{Q}}
\newcommand{\R}{\mathbb{R}}
% \newcommand{\Aut}{\mathrm{Aut}}
% \newcommand{\Tens}{\mathrm{Tens}}
\renewcommand{\d}{\mathrm{d}}
\newcommand{\diff}{\mathrm{d}}
\newcommand{\df}{\Omega}
\newcommand{\Cinfty}{\mathscr{C}^\infty}
\newcommand{\T}{\mathrm{T}}
\newcommand{\cT}{\mathrm{T}^\ast}
\newcommand{\Id}{\mathrm{Id}}
\newcommand{\Lie}{\mathscr{L}}
% \newcommand{\Weil}{\mathcal{W}}
\newcommand{\vf}{\mathfrak{X}}
\newcommand{\X}{\mathfrak{X}}
\newcommand{\Xh}{\X_{\mathrm{ham}}}
\newcommand{\bomega}{\bm{\omega}}

\renewcommand{\L}{\mathcal{L}}
\newcommand{\F}{\mathcal{F}}
\renewcommand{\H}{\mathcal{H}}
\renewcommand{\P}{\mathcal{P}}
\newcommand{\U}{\mathcal{U}}
\newcommand{\W}{\mathcal{W}}
\newcommand{\Reeb}{\mathcal{R}}
\newcommand{\calX}{\mathcal{X}}
\newcommand{\calZ}{\mathcal{Z}}
\newcommand{\C}{\mathcal{C}}
\newcommand{\D}{\mathcal{D}}
\newcommand{\V}{\mathcal{V}}
\renewcommand{\O}{\mathcal{O}}
\renewcommand{\W}{\mathcal{W}}

\newcommand{\g}{\mathfrak{g}}

\newcommand{\Rt}{R_t}
\newcommand{\Rz}{R_z}
\let\Rs\Rz

\newcommand{\rmC}{\mathrm{C}}
\newcommand{\rmR}{\mathrm{R}}
\newcommand{\rmS}{\mathrm{S}}

\newcommand{\bfX}{\mathbf{X}}
\newcommand{\bfY}{\mathbf{Y}}
\newcommand{\bfZ}{\mathbf{Z}}
\newcommand{\bfJ}{\mathbf{J}}

\newcommand{\parder}[2]{\frac{\partial #1}{\partial #2}}
\newcommand{\dparder}[2]{\dfrac{\partial #1}{\partial #2}}
\newcommand{\tparder}[2]{\partial #1/\partial #2}
\newcommand{\parderr}[3]{\frac{\partial^2 #1}{\partial #2\partial #3}}
\newcommand{\dparderr}[3]{\dfrac{\partial^2 #1}{\partial #2\partial #3}}
\newcommand{\tparderr}[3]{\partial^2 #1/\partial #2\partial #3}

\let\graph\relax
\DeclareMathOperator{\graph}{graph}

\DeclareMathOperator{\Ima}{Im}
\DeclareMathOperator{\Ad}{Ad}
\DeclareMathOperator{\ad}{ad}
\DeclareMathOperator{\inn}{\iota}
% \DeclareMathOperator{\coker}{coker}
\DeclareMathOperator{\tr}{tr}
\DeclareMathOperator{\supp}{Supp}
\DeclareMathOperator{\Hom}{Hom}
\DeclareMathOperator{\Ham}{Ham}
% \DeclareMathOperator{\sgn}{sgn}
\DeclareMathOperator{\Diff}{Diff}
% \DeclareMathOperator{\Ad}{Ad}
\DeclareMathOperator{\rk}{rank}
\DeclareMathOperator{\Div}{div}
\DeclareMathAlphabet{\mathpzc}{OT1}{pzc}{m}{it}
\def\d{\mathrm{d}}
\DeclareMathOperator{\id}{id}
\DeclareMathOperator{\pro}{pr}
\newcommand{\pr}[1]{\pro_{#1}}

\newcommand{\liedv}[1]{\Lie_{#1}}
\newcommand*{\contr}[1]{i_{#1}}


% \newtheorem{corollary}{Corollary}
% \newtheorem{definition}{Definition}
% \newtheorem{lemma}{Lemma}\newtheorem{remark}{Remark}
% \newtheorem{proposition}{Proposition}
% %{\theorembodyfont{\rmfamily}\newtheorem{remark}{Remark}}
% \newtheorem{theorem}{Theorem}
%%%%%%%%%% End TeXmacs macros

\newcommand{\comment}[1]{{\noindent}

\begin{flushleft}
  \tmtexttt{#1}
  
  \marginpar{\tmtextsc{Comment}}
\end{flushleft}}
\providecommand{\todo}[1]{{\vspace{5 mm}}

{\noindent}\marginpar{\tmtextsc{ToDo}}

\fbox{\document{\minipage{c}{0.95par}{\concat{ }{#1}}}}

{\vspace{5 mm}}}


\newcommand{\bfi}{}
%\newcommand{\thefigure}{{\thesection}.{\arabic}{\cfigure}}
%\newcommand{\thetable}{{\thesection}.{\arabic}{\ctable}}
\newcommand{\fps}{h, t}
%\newcommand{\theequation}{{\thesection}.\arabic{equation}}


%\newtheorem{examples}{Examples}

\newcommand{\boxit}[1]{{\lrbox{{\savepar}}{ \begin{minipage}[b]{15.8cm}
  #1
\end{minipage}{

}}}\fbox{\usebox{\savepar}}}



\title{{\sffamily Hamiltonian stochastic Lie systems  and applications}}

\author{{\sffamily 
$^a$Javier de Lucas%
\thanks{e-mail:
   javier.de.lucas@fuw.edu.pl \ ORCID: 0000-0001-8643-144X}\ ,\
$^b$Xavier Rivas%
\thanks{e-mail:
   xavier.rivas@unir.net \ ORCID: 0000-0002-4175-5157}, and\
$^c$Marcin Zaj\k ac%
\thanks{e-mail:
   marcin.zajac@fuw.edu.pl \ ORCID: 0000-003-4914-3531}
}
\\[1ex]
\normalsize\itshape\sffamily
$^a$UW Institute for Advanced Studies,\\
\normalsize\itshape\sffamily
Department of Mathematical Methods in Physics, University of Warsaw, Warszawa, Poland.
\\[1ex]
\normalsize\itshape\sffamily
$^b$Escuela Superior de Ingenier\'{\i}a y Tecnolog\'{\i}a,\\
\normalsize\itshape\sffamily
Universidad Internacional de La Rioja, Logro\~no, Spain.
\\[1ex]
\normalsize\itshape\sffamily
$^c$  Department of Mathematical Methods in Physics,
\\
\normalsize\itshape\sffamily
University of Warsaw, Warszawa, Poland.
\\[1ex]
}

%\date{{\sffamily \today}}


\pagestyle{myheadings}
\markright{{\rm J. de Lucas, X. Rivas, and M. Zaj\k ac ---  Hamiltonian stochastic Lie systems and beyond}}



\begin{document}


\maketitle
\begin{abstract}

This paper provides a practical approach to stochastic Lie systems, i.e.  stochastic differential equations whose general solutions can be written as a function depending only on a generic family of particular solutions and some constants, so as to emphasise their applications. We correct the known stochastic Lie theorem characterising stochastic Lie systems, proving that, contrary to previous claims, it satisfies the Malliavin's principle. Meanwhile, we show that stochastic Lie systems admit new stochastic features in the It\^{o} approach. New generalisations of stochastic Lie systems, like the so-called stochastic foliated Lie systems, are devised. Subsequently, we focus on stochastic (foliated) Lie systems that can be studied as Hamiltonian systems using different types of differential geometric structures. We study their stability properties and we devise the basics of an energy-momentum method. A stochastic Poisson coalgebra method is developed to derive superposition rules for Hamiltonian stochastic Lie systems. Applications of our results are found in coronavirus stochastic models, stochastic Lotka--Volterra systems, stochastic SIS models of different types, etc. Our results improve previous approaches by using stochastic differential equations instead of deterministic models designed to grasp some of their stochastic features.
\end{abstract}


\noindent\textbf{Keywords:} 
energy-momentum method,
Hamiltonian stochastic Lie system,
Poisson coalgebra, 
stability,
stochastic differential equation, 
stochastic Lie system,
superposition rule,
\bigskip

\noindent\textit{{\bf MSC 2020:} 
60H10, %Stochastic ordinary differential equations
%60H30, %Aplications of stochastic analysis,
34A26 % Geometric methods in differential equations
(Primary)
%34F05, 
37N25, %Dynamical systems in biology
53Z05. % Symplectic/contact geometry - Applications to physics
(Secondary)
}

\bigskip

{\setcounter{tocdepth}{2}
\def\baselinestretch{1}
\small
% hack per a eliminar l'espai vertical 1em
\def\addvspace#1{\vskip 1pt}
\parskip 0pt plus 0.1mm
\tableofcontents
}

\section{Introduction}

In its most classical version, a {\it superposition rule} is a manner to describe the general solution of a time-dependent system of first-order ordinary differential equations (ODEs) in normal form, a so-called {\it Lie system}, via a time-independent function of a generic set of particular solutions and a set of constants \cite{BM10,CGM00,CGM07,PW83}. Superposition rules have applications, for instance, in the theory of approximate and numerical methods, as they are available for Lie systems whose exact general solutions in explicit form are unknown \cite{Pi12,PW83}. In particular, the knowledge of a particular finite set of exact and/or approximate solutions of a Lie system permits one to study its general solution via superposition rules.

The superposition rule concept traces its origins back to the Sophus Lie's pioneering and celebrated book \cite{Lie1893}. In that work, Lie stated the nowadays called Lie theorem characterising Lie systems. Moreover, \cite{Lie1893} laid down the foundations for the theory of Lie systems. The Lie theorem is also called Lie--Scheffers theorem, as Scheffers took part in writing \cite{Lie1893}, or Lie superposition theorem \cite{BM10local-global}. As  shown in \cite{LS21} and references therein, Lie, alone, described his Lie theorem in \cite{Li93a} without a proof as a criticism to some previous works by Vessiot, Guldberg, K\"oningsberger, and others researchers (see \cite{Gu93,LS21,Ve93} and references therein). Lie stated that the results on superposition rules for differential equations on $\mathbb{R}$ devised by previous authors were a simple application of his theory on infinitesimal group transformations. This suggests one to call Lie theorem the Lie--Scheffers theorem and  to use the denomination `Lie system' instead of  `Lie--Scheffers system'. 

Since Vessiot devised important contributions to the theory of Lie systems \cite{Dissertationes,LS21},  Lie--Vessiot systems is also a good denomination for Lie systems. On the contrary, Scheffers did not work on Lie systems apart from collaborating in writing Lie's works \cite{LS21}. 

The Lie theorem states that a time-dependent system of ODEs in normal form admits a superposition rule if and only if it can be considered as a curve in a finite-dimensional Lie algebra of vector fields, a {\it Vessiot--Guldberg Lie algebra} (see \cite{BM10local-global,CGM00,CGM07} for modern approaches  and further details on this result). 

Lie systems have been thoroughly studied in physics and mathematics due to their widespread occurrence  (see \cite{Dissertationes,LS21} counting more than 300 references on Lie systems and related topics). %The classical Lie Theorem {} characterises time-dependent systems of first-order differential equations in normal form admitting a superposition rule.
There has been a bast effort to generalise Lie systems to more general situations: time-dependent Schr\"odinger equations \cite{Dissertationes}, quasi--Lie systems \cite{CGL19} and quasi--Lie families \cite{CL15}, foliated Lie systems \cite{CLW22}, etc \cite{CL09}. In this work, we are mainly concerned with the extension of superposition rules and Lie systems to the realm of stochastic differential equations \cite{LO09}. %Although there are generalisations of superposition rules for higher-order  and other types of differential equations, we will be hereafter dealing always with deterministic or stochastic time-dependent first-order ordinary differential equations.

Stochastic differential equations may describe phenomena that  deterministic differential equations can not \cite{Al09}. For example, the probability of disease extinction or outbreak, the quasi-stationary probability distribution, the final size distribution, and the expected duration of an epidemic are features that cannot  effectively be modelled by deterministic methods. These and other reasons motivate the huge interest in studying stochastic differential equations. Although some deterministic differential equations can grasp certain characteristics of stochastic differential equations \cite{CFH23}, only strict stochastic differential equations offer a complete view on the above-mentioned phenomena. In particular, it would be interesting to extend to stochastic differential equations the concept of superposition rule. 

The work \cite{LO09} extends the superposition rule notion to a class of systems of stochastic first-order ordinary differential equations, called  {\it stochastic Lie systems}. Moreover, \cite{LO09} presents a nice and very precise presentation of certain local results about stochastic Lie systems and, as a byproduct, it also explains many technical results on standard Lie systems. Despite its mathematical interest, such details are usually skipped in the literature as they have, in general, not many practical applications. As noted in \cite{BM10local-global,LS21}, standard works on Lie systems, even theoretically ones, are essentially interested in local aspects and generic points, which makes them skip many technical details analysed in \cite{LO09}. Notwithstanding, there is mathematical interest in global and local technical aspects, as illustrated by the work \cite{BM10local-global} for global superposition rules and \cite{LO09}. It is worth stressing that stochastic Lie systems were found to have applications in the description of Brownian motions, economical models like the Black-Scholes theory of derivatives prices, and so on \cite{LO09}.

In \cite{LO09}, a stochastic Lie theorem  characterising Stratonovich stochastic first-order ordinary differential equations admitting a superposition rule was devised, but it contains a mistake that  changes its meaning and potential applications. More exactly, the direct part of the stochastic Lie theorem in \cite{LO09} states that a stochastic Lie system may admit, locally, a superposition rule if its associated Stratonovich operator is related to a family of vector fields spanning an involutive distribution. Our present work proves that the vector fields must additionally span a finite-dimensional real Lie algebra of vector fields. We also determine the exact point of the mistake in \cite{LO09}, we correct the enunciate of the stochastic Lie theorem, and we provide a counterexample to illustrate when the stochastic Lie theorem in \cite{LO09} fails. In fact, we prove that this mistake  leads to characterise certain stochastic Lie systems related to SIS epidemic models, when they are not. 
On the other hand, the differences between the wrong and the correct versions of the stochastic Lie theorem are assumed to occur in a case that does not occur in most accomplished applications in \cite{LO09}.% In fact, as far as we known, there is just one application in \cite{LO09}, where the stochastic Lie theorem was wrongly applied. 

Moreover, our work presents a concise introduction to stochastic Lie systems, aiming at providing a practical approach while avoiding  technical details that are not necessary for our purposes. In this sense, it simplifies the elegant and mostly rigorous mathematical treatment in  \cite{LO09} by using standard assumptions in mathematical constructions. For instance, we focus on local results at generic points, which significantly simplifies previous needed techniques. 

We here show that the theory of stochastic Lie systems in the Stratonovich approach reassembles exactly what one could expect from the standard theory of Lie systems. Indeed, our correction of the stochastic Lie theorem shows that, contrary to what it was claimed in \cite{LO09}, the stochastic Lie theorem has no proper stochastic features in the Stratonovich approach. Meanwhile, it is stressed that the conditions for a stochastic differential equation in the It\^{o} form \cite{Ka81} to become a stochastic Lie system do not reassemble the standard form expected from the deterministic Lie theorem \cite{CGM00,CGM07,Dissertationes}. This is very important practically, as many relevant stochastic differential equations are given in the It\^o form and they have to be translated into the Stratonovich approach \cite{Sa70} to apply the methods of our work and \cite{LO09}. In this respect, the relation between the It\^o and the Stratonovich approaches is reviewed and some examples are provided in this work. Special attention is paid to show that stochastic Lie systems take an unexpected form in the It\^{o} form, which does not follow the standard classical form in the deterministic theory. %This is specially important as the Stratonovich and It\^{o} approaches frequently in applications \cite{Sa70}.  

Apart from introducing stochastic Lie systems, this paper suggests the applicability of various generalisations of Lie systems \cite{Dissertationes} to the stochastic domain. This unveils a vast realm of potential applications, offering promising avenues for further exploration. One notable example is the extension of the theory of foliated Lie systems \cite{CLW22} to the realm of stochastic differential equations. In fact, we suggest that, in analogy with the deterministic case, this can be the case when  studying certain problems of our here devised stochastic energy-momentum method \cite{LZ21}. 
%----------------------------------------------------------

By the stochastic Lie theorem, every stochastic Lie theorem admitting $\ell$ independent random variables is related to a family of $\mathbb{R}^\ell$ vector fields. All of them can be understood as linear combinations of elements of a Vessiot--Guldberg Lie algebra of vector fields. We study stochastic Lie systems admitting a Vessiot--Guldberg Lie algebra of Hamiltonian vector fields relative to some compatible differential geometric structure: the {\it Hamiltonian stochastic Lie systems}. In particular, we focus on Hamiltonian stochastic Lie systems relative to symplectic structures. In this context, the coalgebra method is extended to Hamiltonian stochastic Lie systems to derive superposition rules. This provides an extension to the stochastic realm of the theory of Hamiltonian Lie systems and their generalisations \cite{LS21}. 

Our results are illustrated with many new examples. In particular, we study several stochastic new and known {\it SIS systems} \cite{CCM19,CCM21,EFSZ22,NM19,Ot21,WCDG18}. SIS models are epidemiological models that assume that individuals do not acquire immunity after infection. They concern two variables/compartments: $S$ representing susceptible individuals and $I$ representing infected individuals in a large population of size $N$ where a single disease is spreading. Similarly, SIS models can also describe the spreading of gossips over the internet and other similar effects.  In the SIS model, the most important variable is the instantaneous density of infected individuals,  $\rho = \rho(\tau)$, which depends on the time parameter $\tau$ and ranges between 0 and 1.
The density of infected individuals decreases at a rate of $\gamma\rho$, where $\gamma$ is the recovery rate. The rate of new infections is proportional to $\alpha\rho(1-\rho)$, where $\alpha$ represents the transmission rate or intensity of contagion. These two processes are described by the following compartmental equation:
\[
\frac{\d\rho}{\d\tau} = \alpha\rho(1-\rho) - \gamma\rho\,.
\]
To simplify the equation, define  $t := \alpha\tau$, for the case $\alpha>0$, and introduce a constant $\rho_0 := 1 - \dfrac{\gamma}{\alpha}$. Then,

\[
\frac{\d\rho}{\d t} = \rho(\rho_0-\rho)\,.
\]
This formulation allows for a clearer representation of the dynamics of the SIS model.

Such systems are usually employed in the literature in a deterministic manner. This can be used to describe some of their features, but not all, as some are pure stochastic. In particular, we try to provide a more general approach than in \cite{CFH23,EFSZ22}, where SIS systems are studied in a deterministic manner without using stochastic methods. Our models can be also used to study stochastic models appearing in Lotka--Volterra models. Moreover, the systems in \cite{CFH23} can also be generalised to a real stochastic realm and be described via Hamiltonian stochastic Lie systems. 

%The notion is motivated from examples present in epidemiological models. In particular, we study stochastic Lie systems that SIR systems via symplectic and contact structures. 

We are also concerned with a theory of stability for Hamiltonian stochastic systems and, in particular, Hamiltonian stochastic Lie systems. Our study is specially concerned with linear ones, which appear as an approximation of nonlinear ones and they share some stability properties with them \cite{Ar74}. Moreover, the basis for a energy-momentum method \cite{MS88} for Hamiltonian stochastic differential equations is laid down by using some results in \cite{LO08}. In particular, this allows one to study the relative equilibrium points of Hamiltonian systems, i.e. points where the dynamics is generated by Hamiltonian symmetries of the system under study. In particular, a characterisation (see Theorem \ref{RelativeEquilibrum}) of the relative equilibrium points for stochastic Hamiltonian systems in terms of critical points of its Hamiltonian functions is developed.

The structure of the paper goes as follows. 
Section~\ref{Fundamentals} is a brief introduction to stochastic differential equations, stochastic Lie systems, and many other notions to be used in this paper. It also shows the difference of the form of stochastic Lie systems in the Stratonovich and the It\^{o} approaches. Section~\ref{Sec:SSRandLT} is concerned with superposition rules for stochastic Lie systems and it reviews and corrects the previous version of the stochastic Lie system. Section \ref{Sec:Hamiltonian_Stochastic_Lie_Systems} deals with Hamiltonian stochastic Lie systems. In particular, we define the newly proposed stochastic foliated Lie systems, and the Hamiltonian counterparts of stochastic Lie systems. A theory of stability for stochastic Lie systems is given in Section \ref{Sec:SHS}. A relative equilibrium notion is presented and studied in Section \ref{Sec:Relative_Equilibrium_Stochastic_Differential}. Meanwhile, Section \ref{Se:PCSLS} develops the Poisson coalgebra method for Hamiltonian stochastic Lie systems. Finally, Section \ref{Eq:Conclusions} presents our conclusions and future work.  

\section{Stochastic differential
equations and stochastic Lie systems}\label{Fundamentals}

This section provides a concise introduction to stochastic differential equations and stochastic Lie systems. The detailed survey on the theory of stochastic differential equations can be found in \cite{Ar74, E82}, while the theory of the stochastic Lie systems was elaborated for the first time in \cite{LO09}, which offers a precise theory. To focus on the main parts of the theory, we will assume objects to be smooth and locally defined. Following the classical Lie theory \cite{Lie1893}, we here provide a stochastic Lie system definition not based on the superposition rule notion.
%-------------------------------------------------------

In a nutshell, a time-dependent system of first-order differential equations on an $n$-dimensional manifold $M$ of the form
\begin{equation}\label{eq:Det}
\frac{\d \Gamma^i}{\d t}=X^i(t,\Gamma)\,,\qquad i=1,\ldots,n\,,\qquad \forall t\in \mathbb{R}\,,\qquad \forall \Gamma\in M\,,
\end{equation}
for certain functions $X^1,\dotsc,X^n\in \Cinfty(\mathbb{R}\times M)$,
 is deterministic in the sense that an initial condition in $M$ establishes, under mild conditions on  $X^1,\ldots,X^n$, a unique solution giving a position in $M$ for every $t\in \mathbb{R}$. Then, \eqref{eq:Det} may be modified by considering that its solutions can depend on certain `stochastic processes' $B^1,\ldots,B^r$, namely a series of $t$-dependent random variables. This fact is shown by the expression
\begin{equation}\label{Eq:practicalSDE}
\delta \Gamma^i=X^i(t,B,\Gamma)\delta t+\sum_{\alpha=1}^rX^i_\alpha(t,B,\Gamma)\circ \delta B^\alpha\,,\qquad i=1,\ldots,n\,,
\end{equation}
for functions $X^i , X^i_1 ,\ldots,X_r^i $, with $B=(B^1,\ldots,B^r)$, $\Gamma=(\Gamma^1,\ldots,\Gamma^n)$, and $i=1,\ldots,n$. The symbol $\circ$ has been used to indicate that \eqref{Eq:practicalSDE} is understood in the so-called Stratonovich interpretation, to be briefly explained afterwords.
More formally and precisely, $(\Omega, \mathcal{F}, P)$ is a {\it probability space}, where $\Omega$ is a manifold, $\mathcal{F}$ is a $\sigma$-algebra of subsets of $\Omega$, and $P:\mathcal{F}\rightarrow [0,1]$ is a probability function on $\mathcal{F}$. Each $B_\alpha:\mathbb{R}_+\times\Omega\rightarrow \mathbb{R}$ is a semi-martingale for $\alpha=1,\ldots,r$. Semi-martingales are good integrators relative to the It\^{o} or the Stratonovich integrals due to its properties. %In particular, a semi-martingale, $X:\mathbb{R}_+\times \Omega\rightarrow \mathbb{R}$ is the sum $X_t=M_t+A_t$ of a martingale $M_t$ and a c\`adl\`ag\footnote{from the french ``continue à droite, limite à gauche'', namely ``right-continuous with left limits''.} adapted process of locally bounded variation $A_t$. 
A martingale  is a sequence of stochastic processes such that, at a particular time, the conditional expectation of the next value is equal to the present value, independently of all previous values.

A stochastic differential equation is then an expression on a manifold $M$ of the form
\begin{equation}
  \label{stochastic differential equation expression} \delta \Gamma = \mathfrak{S} (X,
  \Gamma) \delta X\,,
\end{equation}
where $X:\mathbb{R}_+ \times \Omega \rightarrow \R^\ell$ is an
$\R^\ell$-valued semi-martingale and $\mathfrak{S} (X, \Gamma) : \T_X \mathbb{R}^\ell
\to \T_\Gamma M$, where $(X,\Gamma)\in \mathbb{R}^\ell\times M$, describes a Stratonovich operator. Every basis in
$\T^{\ast} \R^\ell$ identifies $\mathfrak{S}(X,
\Gamma)$ as an $\ell$-tuple of its components $(\mathfrak{S}_1 (X, \Gamma), \ldots, \mathfrak{S}_\ell (X, \Gamma))$ in the chosen basis. Every particular solution to \eqref{stochastic differential equation expression} is also a semi-martingale $
\Gamma:\mathbb{R}_+\times \Omega\rightarrow M$. Moreover, we say that a particular solution has initial condition $\Gamma_0
\in M$ when $\Gamma(0,\omega)=\Gamma_0$ for every $\omega\in \Omega$ with probability one. Note that the standard time can be considered as a random variables $t:(t,\omega)\in \mathbb{R}_+\times \Omega\mapsto t\in\mathbb{R}$ whose value is independent of $\Omega$ and to be included as the first component of $X$.

The solution to \eqref{stochastic differential equation expression} is considered to be of the form
\begin{equation}\label{Eq:SolStr}
\Gamma(t)-\Gamma(0)=\int_0^t\mathfrak{S}_1(X,\Gamma)\delta t+\sum_{\beta=2}^{\ell}
\int_0^t\mathfrak{S}_\beta(X,\Gamma) \circ \delta X^\beta_t,
\end{equation}
where the integrals appearing above are {\it Stratonovich integrals}. It is relevant to understand that stochastic differential equations can still be understood in the so-called It\^{o} way, in which the general solution is also of the form \eqref{Eq:SolStr}, but the integrals are assumed to be in the It\^{o} manner, which is different. If one approach or the other is used, depends on the applications to be developed.

Unless otherwise stated, every stochastic differential equation \eqref{stochastic differential equation expression} is assumed to be in the Stratonovich way. This is motivated by the Malliavin’s
Transfer Principle \cite{Ma78}, we suggest, but not prove, that the obtained theory will reassemble the standard Lie theory. Despite this general idea, stochastic differential equations appear in the It\^{o} sense in many applications. Hence, a manner to deal with such stochastic differential equations will be studied in this work. Moreover, the study of such stochastic differential equations show some differences appearing in the stochastic formalism with respect to the deterministic theory of Lie systems, which much enriches the theory.

Although every Stratonovich stochastic differential equation is equivalent to another It\^{o} stochastic differential equation, the form of both is different. Hence, one has to take care of the method is employed to study a stochastic differential equation \cite{Ev13}. More exactly, assume that \eqref{Eq:practicalSDE} has coefficients that do not depend on $\mathcal{B}^2,\ldots,\mathcal{B}^\ell$, then the It\^{o} differential equation admitting the same solutions  reads (see \cite{BCC_01} for details)
\begin{equation}\label{Eq:TranItoStratonovich}
\delta \Gamma^i=\left(\mathfrak{S}_1(t,\Gamma)+\frac 12\sum_{\beta=2}^\ell\sum_{j=1}^n\frac{\partial \mathfrak{S}_\beta^i}{\partial \Gamma^j}(t,\Gamma)\mathfrak{S}_\beta^j(t,\Gamma)\right)\delta t +\sum_{\beta=2}^\ell \mathfrak{S}^i_\beta(t,\Gamma )\delta\mathcal{B}^\beta_t,\qquad i=1,\ldots,n.
\end{equation}
Note that we have dropped the $\circ$ sign in the previous expression.
It is worth stressing that the definition of a stochastic Lie system to be given soon relies on the new term appearing multiplying $\delta t$ in \eqref{Eq:TranItoStratonovich}, which is called the {\it drift term} \cite{BCC_01}. Moreover, the transformation from the Stratonovich to the Itô form does not change the coefficients with the $\delta\mathcal{B}^\alpha_t$ for $\beta = 2,\dotsc,\ell$.

\begin{example}\label{Ex:FirstItoStr} Let us consider a stochastic differential equation induced by a semi-martingale $X:(t,z)\in \mathbb{R}_+\times \Omega \mapsto (t,\mathcal{B})\in \mathbb{R}^2$ consisting of two variables (one deterministic, $t$, which can be understood as a stochastic variable, and one purely stochastic, $\mathcal{B}$) of the It\^{o} form
\begin{equation}\label{eq:It}
\delta I=I(\beta N-\mu-\gamma-\beta I)\delta t+I\sigma (N-I) \delta \mathcal{B}\,,
\end{equation}
where $N$ is a constant describing the total population of a SIS epidemiological system \cite{GGHMP11}. Let us assume that $\sigma=\sigma(t)$. The Stratonovich stochastic differential equation related to \eqref{eq:It} takes the form
\begin{equation}\label{eq:St}
\delta I=\left(-\sigma^2(t) I^3+\left(-\beta+\frac{3N\sigma^2(t)}{2}\right)I^2+I\left(N \beta-\gamma-\mu-N^2\frac{\sigma^2(t)N^2}{2}\right)\right)\delta t+\sigma(t) I(N-I)\circ \delta \mathcal{B}\,.
\end{equation}
This induces a Stratonovich operator between $\mathbb{R}^2$ and $\mathbb{R}$ such that, for each $t,\mathcal{B},I$, one has an operator
$$ \mathfrak{S}(t,\mathcal{B},I):(\delta t,\delta \mathcal{B})\in \T_{(t,\mathcal{B})}\mathbb{R}^2\longmapsto \mathfrak{S}_1(t,I)\delta t+\mathfrak{S}_1(t,I)\delta \mathcal{B}\in \T_I\mathbb{R}\, $$
for 
$$
\begin{gathered}
\mathfrak{S}_1(t,I)=-\sigma^2(t) I^3+\left(-\beta+\frac{3N\sigma^2(t)}{2}\right)I^2+I\left(N \beta-\gamma-\mu-N^2\frac{\sigma^2(t)N^2}{2}\right),\\
\mathfrak{S}_2(t,I)=\sigma(t) I(N-I).
\end{gathered}
$$
Solutions to \eqref{eq:St} are given by semi-martingales of the form $\Gamma:(t,\omega)\in \mathbb{R}_+\times \Omega\mapsto I(t,\omega)\in \mathbb{R}$ determined by an initial condition described by a random variable $\Gamma_0:\Omega\rightarrow \mathbb{R}$ such that $
\Gamma_0(\omega)=1$ for every $\omega\in \Omega$. \demo
\end{example}

The theoretical utilisation of Stratonovich stochastic equations is justified by Malliavin's Transfer Principle \cite{Ma78}, which states that the results from the theory of ordinary differential equations remain applicable, in an analogous way, to stochastic differential equations in Stratonovich form. This principle is just a general conception without a proof, which implies that it must be used just as a general indication. Nevertheless, as shown in Example \ref{Ex:FirstItoStr}, the relation to the It\^{o} approach has to be considered too for analysing applications.

\begin{definition} A {\it stochastic Lie system} is a stochastic differential equation on a manifold $M$ of the form $
\delta \Gamma=\mathfrak{S}(X,\Gamma)\circ \delta X
$ such that $X:\mathbb{R}_+\times \Omega\rightarrow \mathbb{R}^\ell$ is a semi-martingale and $\mathfrak{S}$ is a Stratonovich operator such that 
\begin{equation}
\label{Eq:STL}\mathfrak{S}(X,\Gamma)=\left(\sum_{\alpha=1}^rb_1^\alpha(X)Y_\alpha(\Gamma),\ldots,\sum_{\alpha=1}^rb_\ell^\alpha(X)Y_\alpha(\Gamma)\right),\qquad \forall\Gamma\in M,\quad \forall X\in \mathbb{R}^\ell,
\end{equation}
for family of functions $b^a_\alpha:x\in X\mapsto b^a_\alpha(X)\in \mathbb{R}$, with $\alpha=1,\ldots,r$ and $a=1,\ldots,\ell$, a certain $r$-dimensional Lie algebra of vector fields spanned by $Y_1,\ldots,Y_r$.    We call the Lie algebra $\langle Y_1,\ldots,Y_r\rangle$ a {\it Vessiot--Guldberg Lie algebra} of the stochastic Lie system \eqref{Eq:STL}.
\end{definition}

Note that the Stratonovich operator $\mathfrak{S}$ of a stochastic Lie system can be considered as a series of $\ell$-vector fields depending on $X$. Each of its components can be understood as a section of the trivial bundle $ \mathbb{R}^{\ell}\times V\rightarrow \mathbb{R}^\ell$, and conversely.




\begin{example}\label{eq:WhiteOscil} Let us consider a damped harmonic oscillator on $\T \mathbb{R}$ with adapted coordinates $\{x,y\}$ and with white noise $W_1$ (see \cite{Ha19}) given by
\begin{equation}\label{eq:HOWN}
\delta\begin{pmatrix} x \\ y \end{pmatrix} = \begin{pmatrix} 0 & \omega \\ -\omega & -k \end{pmatrix}\begin{pmatrix} x \\ y \end{pmatrix}\delta t + \begin{pmatrix} 0 & 0 \\ 0 & -\sigma \end{pmatrix}\begin{pmatrix} x \\ y \end{pmatrix} \delta W_1\,.
\end{equation}
More specifically, $W_1$ is a one-dimensional Wiener process, while $\sigma$ is a parameter quantifying the noise. The parameters $\omega$ and $k$ are the usual parameters relative to the standard deterministic model for a dissipative harmonic oscillator
$$ \ddot x+\omega^2 x+k\dot x=0\,, $$
where $\omega^2$ is the elastic constant of the oscillator and $k$ the friction coefficient. Let us see how this model can be considered as a linear stochastic Lie system.

Consider the vector fields on $\T \mathbb{R}$ given by
\begin{equation}\label{Eq:BaGL22}
 X_{11}=x\frac{\partial}{\partial x}\,,\qquad X_{12}=y\frac{\partial}{\partial x}\,,\qquad X_{21}=x\frac{\partial}{\partial y}\,,\qquad X_{22}=y\frac{\partial}{\partial y}\,.
\end{equation}
 These vector fields span a Lie algebra of vector fields $V_D$ isomorphic to the general Lie algebra, $\mathfrak{gl}(2,\mathbb{R})$, of $2\times 2$ matrices with real coefficients. In fact, the commutation relations read
 \begin{align*}
 & [X_{11},X_{12}] = -X_{12}\,, && [X_{11}, X_{21}] = X_{21}\,, && [X_{11}, X_{22}] = 0\,,\\
 & [X_{12}, X_{21}] = X_{22} - X_{11}\,, && [X_{12},X_{22}] = -X_{12}\,,\\
 & [X_{21},X_{22}] = X_{21}\,.
 \end{align*}
Moreover, \eqref{eq:HOWN} can be related to a Stratonovich operator of the form
$$ \mathfrak{S}(t, W_1, x, y)=\left(\begin{pmatrix} 0 & \omega \\ -\omega & -k \end{pmatrix}\begin{pmatrix}x\\y\end{pmatrix},\begin{pmatrix}0 & 0 \\ 0 & -\sigma\end{pmatrix}\begin{pmatrix}x\\y\end{pmatrix}\right). $$
Hence, one can write
$$
\mathfrak{S}(t, W_1, x, y)=\left(\omega (X_{12}-X_{21})-kX_{22},-\sigma X_{22}\right). 
$$
which turns \eqref{eq:HOWN} into a stochastic Lie system and $V_D$ into an associated Vessiot--Guldberg Lie algebra. Note that more general models can be considered by assuming the coefficients $\omega$, $\sigma$, or $k$ depending on $t$ and/or $W_1$. These models will be still stochastic Lie systems.
\demo
\end{example}

\begin{example} It is worth noting that an It\^{o} stochastic differential equation may seem to have the form \eqref{Eq:STL} without being a stochastic Lie system. This shows that the stochastic case does not matches exactly the form given in classical Lie systems. Let us provide a counter example of this. The SIS model
describes the dissemination of a single communicable disease in a susceptible population of size $N$ (see \cite{GGHMP11}).
Let us study a stochastic differential equation that is related to the deterministic SIS model for a particular value of $N$, which is indeed a deterministic approximation of it. Consider the It\^{o} stochastic system
\begin{equation}\label{Eq:PseudoSto}
\delta I=I(5-I/2)\delta t+I\sigma(100-I)\delta B\,,
\end{equation}
for a certain time-dependent parameter $\sigma(t)$, which generalises the stochastic SIS system studied in \cite{GGHMP11}. Consider the vector fields
 $$
 Y_1=\frac{\partial}{\partial I},\qquad Y_2=I\frac{\partial}{\partial I},\qquad Y_3=I^2\frac{\partial}{\partial I},
 $$
 which span a Lie algebra of vector fields with commutation constants
 $$
 [Y_1,Y_2]=Y_2,\qquad [Y_1,Y_3]=2Y_2,\qquad [Y_2,Y_3]=Y_3.
 $$
 Moreover, the operator of the form
 $$
 \mathcal{I}(t,B,I)=(5Y_2-1/2Y_3,100\sigma(t) Y_2-\sigma(t) Y_3),
 $$
is related to the Vessiot--Guldberg Lie algebra $\langle Y_1,Y_2,Y_3\rangle$ and one would expected \eqref{Eq:PseudoSto} to be a stochastic Lie system. Nevertheless, one has that \eqref{Eq:PseudoSto} is related to the Stratonovich differential equation
\begin{equation}\label{Eq:StrSIS}
\delta I=\left(-\sigma(t)I^3+\left(-\frac 12+150\sigma(t)^2\right)I^2+I(5-5000\sigma^2(t))\right)\delta t+I\sigma(t)(100-I)\circ \delta B\,,
\end{equation}
which cannot be described by Vessiot--Guldberg Lie algebra $\langle Y_1,Y_2,Y_3\rangle$. Indeed, one can prove, as shown hereafter in Section \ref{Sec:SSRandLT}, that \eqref{Eq:StrSIS} is not a stochastic Lie system.
\demo
\end{example}


There are  It\^{o} stochastic differential equations whose coefficients match the form of the coefficients in \eqref{Eq:STL} and they are still stochastic Lie systems. This is due to the fact that the transformation \eqref{Eq:TranItoStratonovich} maps to the initial \eqref{Eq:practicalSDE} into a new stochastic differential equation that reassembles again the condition \eqref{Eq:STL} invariant. Notwithstanding, this is not the general case.
\begin{example} Let us consider a satellite in a circular orbit under the influence of a rapidly fluctuating density of the atmosphere of the orbit \cite[Eq. (3.25)]{Sa70}
$$ \ddot Y_t+B(1+A\xi_t)\dot Y_t+(D-2C+AD\xi_t)Y_t=0\,, $$
where $B,C$ are positive numbers and $\xi$ is a scalar white noise. Let us assume this problem to be approximated through a stochastic differential equation \cite[Eq. (3.29)]{Sa70}
\begin{equation}\label{Eq:linAp}
 \delta \Gamma=\delta \begin{pmatrix}
    Y_t\\ \dot Y_t
\end{pmatrix} = \begin{pmatrix}
    0 & 1 \\ 2C-D & -B
\end{pmatrix}\begin{pmatrix}
    Y_t \\ \dot Y_t
\end{pmatrix} \delta t + \begin{pmatrix} 
    0 & 0 \\ -AD & -AB
\end{pmatrix} \begin{pmatrix}
    Y_t\\ \dot Y_t
\end{pmatrix}\delta W_t\,.
\end{equation}
It is worth noting that \cite{Sa70} interprets this stochastic differential equation in the It\^{o} way, which means that one needs to transform it into its Stratonovich version so as to apply our techniques. In particular, one has that the coefficients of the stochastic differential equation are given by
$$
 \mathfrak{I}(t,W,Y,\dot Y)=\left(X_{12}+(2C-D)X_{21}-BX_{22},-ADX_{21}-ABX_{22}\right),
 $$
  in terms of the basis of vector fields on $\T \mathbb{R}$ given by \eqref{Eq:BaGL22} and substituting $x$ by $Y$ and $y$ by $\dot Y$, respectively.
In this case, the Stratonovich operator related to this problem takes the form
 $$
 \mathfrak{S}(t,W,Y,\dot Y)=\left( X_{12}+\left(2C-D-\frac 12A^2 D B\right)X_{21}-\left(B+\frac 12 A^2B^2\right)X_{22},-ADX_{21}-ABX_{22}\right).
 $$
 Hence, \eqref{Eq:linAp} becomes a stochastic Lie system related to a Vessiot--Guldberg Lie algebra $V_D$.  
 \demo
 \end{example}


The following result will be of utility for applications.

\begin{definition} We call {\it noise-free stochastic Riccati differential equation}  the Stratonovich stochastic differential equation on $\mathbb{R}$ with stochastic variables $\mathcal{B}_t^2,\ldots,\mathcal{B}_t^\ell:\mathbb{R}_+\times \Omega\rightarrow \mathbb{R}$ of the form 
\begin{equation}\label{eq:NF_Stochastic_Riccati_Equation}
\delta \Gamma=\left(\sum_{\alpha=0}^2b_\alpha(t)\Gamma^\alpha\right)\delta t+\sum_{\beta=2}^\ell\sum_{\alpha=0}^2 b_{\beta\alpha}(t) \Gamma^\alpha\circ \delta \mathcal{B}^\beta_t,\qquad \forall \Gamma\in \mathbb{R},\qquad \forall t\in \mathbb{R}.
\end{equation}
for arbitrary $t$-dependent functions $b_\alpha(t),b_{\alpha\mu}(t)$ with $\alpha=1,2,3$ and $\mu=2,\ldots,\ell$. 
\end{definition}
Noise-free stochastic Riccati differential equations are stochastic Lie systems. The following proposition is immediate.
\begin{proposition}
An It\^{o} differential equation of the form 
$$
\delta \Gamma=\left(\sum_{\alpha=0}^2b_\alpha(t)\Gamma^\alpha\right)\delta t+\sum_{\beta=2}^\ell\sum_{\alpha=0}^1 b_{\beta\alpha}(t) \Gamma^\alpha \delta \mathcal{B}^\beta_t,\qquad \forall \Gamma\in \mathbb{R},\qquad \forall t\in \mathbb{R}.
$$
is also a noise-free stochastic Riccati differential equation.    
\end{proposition}

 It is worth noting that the theory of Lie systems can be generalised to many different realms \cite{Dissertationes,LS21}. In particular, there is the theory of foliated Lie systems. This suggests us the following generalisations.


\begin{definition}[Foliated stochastic Lie system]
    A {\it stochastic foliated Lie system} is a stochastic system of differential equations on $M$ of the form
  \begin{equation}
    \delta \Gamma = \mathfrak{S}(X, \Gamma) \delta X\,, \label{eq1bis2}
  \end{equation}
  where $X:\R_+ \times \Omega \rightarrow \mathbb{R}^\ell$ is an 
  $\mathbb{R}^\ell$-valued semi-martingale and $\mathfrak{S}(X, \Gamma) : \T_X \mathbb{R}^\ell
  \to \T_\Gamma M$ is a Stratonovich operator, such that 
\begin{equation}\label{decomposition with the Ys}
    \mathfrak{S}_j (X, \Gamma) = \sum_{\alpha = 1}^r b_j^\alpha (X,\Gamma) Y_\alpha (\Gamma)\,,\qquad j=1,\ldots,\ell\,,\qquad \forall X\in \mathbb{R}^\ell,\quad \forall \Gamma\in M,
    \end{equation}
and the vector fields $\{ Y_1, \ldots, Y_r \}$ on $M$ span an $r$-dimensional Vessiot--Guldberg Lie algebra such that
$b^\alpha_j(X)$, with $\alpha=1,\ldots,r$ and $j=1,\ldots,\ell$, are first integrals of the $Y_1,\ldots,Y_r$ for every $X\in \mathbb{R}^\ell$. 
\end{definition}

It is worth noting that there are many potential applications concerning the generalisation to the stochastic realm of famous types of Lie systems such as matrix, projective Riccati equations, Bernouilli equations, and so on with applications in predator-prey models, satellite motion, etcetera \cite{Ar03,Sa70}.  Moreover, nonlinear stochastic differential equations are difficult to study. Notwithstanding, under certain conditions, its linearisation can describe its equilibrium properties \cite{Ar03}. Linear or even affine stochastic differential equations are stochastic Lie systems (cf. \cite{LO09}).

\section{Superposition rules and stochastic Lie systems}\label{Sec:SSRandLT}

Let us study the superposition rule notion for stochastic differential equations and the characterisation of stochastic differential equations admitting a superposition rule. This will lead to review and slightly correct some mistakes in the main theorem in \cite{LO09}. 

\begin{definition}\label{def 1}
    A {\it superposition rule} for a Stratonovich stochastic differential equation of the form~\eqref{stochastic differential equation expression} on a manifold $M$ is a function $\Phi: M^{m + 1} \to M$, such that, for a generic set $\Gamma_1, \ldots, \Gamma_m:\mathbb{R}_+\times \Omega\rightarrow M$, of particular solutions of~\eqref{stochastic differential equation expression}, the general solution $\Gamma$ to~\eqref{stochastic differential equation expression} reads
    \[ \Gamma = \Phi (z
    ; \Gamma_1, \ldots, \Gamma_m), \]
    where $z\in M$ is a point to be related to initial conditions. %A {\it stochastic Lie system} is a stochastic system of differential equations admitting a superposition rule.
\end{definition}

It is remarkable that superposition rules for stochastic differential equations do not depend on $X$ or the probability space. 

Let us introduce now several concepts that will be useful to describe and calculate superposition rules for stochastic Lie systems.

The {\it diagonal prolongation to $M^m$ of a vector bundle} $\tau:F\to M$  is the vector bundle $\tau^{[m]}:F^m = F\times\overset{(m)}{\dotsb}\times F\mapsto M^m= M\times \overset{(m)}{\dotsb} \times M$, of the form
\begin{equation*}
   \tau^{[m]}(f_{(1)},\dotsc,f_{(m)}) = (\tau(f_{(1)}),\dotsc,\tau(f_{(m)}))\,,\qquad \forall f_{(1)},\ldots,f_{(m)}\in F\,, 
\end{equation*}
with
\begin{equation}\label{Eq:Fiber}    F^m_{(x_{(1)},\ldots,x_{(m)})}=F_{x_{(1)}}\oplus\cdots\oplus F_{x_{(m)}}\,,\qquad \forall (x_{(1)},\ldots,x_{(m)})\in M^m\,.
\end{equation}
%If it is clear which is the projection of the vector bundle $\tau^{[m]}$, we will sometimes denote it simply by $E^{[m]}$.

Every section $e:M\to F$ of the vector bundle $\tau$ has a natural {\it  diagonal prolongation} to a section $e^{[k]}$ of the vector bundle $\tau^{[k]}$ given by
\begin{equation*}
    e^{[k]}(x_{(1)},\ldots,x_{(k)})=(e(x_{(1)}),\cdots ,e(x_{(k)}))\,,\qquad \forall (x_{(1)},\ldots,x_{(k)})\in M^k.
\end{equation*}
If we consider that every $e(x_{(a)})$ takes values in the $a$-th copy of $F$ within \eqref{Eq:Fiber}, one can write
\begin{equation*}     e^{[k]}(x_{(1)},\ldots,x_{(k)})=e(x_{(1)})+\cdots +e(x_{(k)}))\,,\qquad \forall (x_{(1)},\ldots,x_{(k)})\in M^k,
\end{equation*}
which is a simple useful notation for applications. The {\it  diagonal prolongation of a function $f\in \Cinfty (M)$} to $M^k$ is the function on $M^k$ given by 
$$f^{[k]}(x_{(1)},\ldots,x_{(k)})= f(x_{(1)})+\ldots+f(x_{(k)})\,.
$$
Consider also the sections $e^{(j)}$ of $\tau^{[k]}$, where $j\in \{1,\ldots,n\}$ and $e$ is a section of $\tau$, given by
\begin{equation}\label{prol1}
e^{(j)}(x_{(1)},\dots,x_{(k)})=0+\cdots +e(x_{(j)})+\cdots+0\,,\qquad \forall (x_{(1)},\ldots,x_{(k)})\in M^k\,.
\end{equation}
If $\{e_1,\ldots, e_r\}$ is a basis of local sections of the vector bundle $\tau$, then $e_i^{(j)}$, with $j = 1,\dotsc,k$ and $i = 1,\dotsc,r$, is a basis of local sections of $\tau^{[k]}$.

Due to the obvious canonical isomorphisms
$$(\T M)^{[k]}\simeq \T M^k\quad\text{and}\quad (\cT M)^{[k]}\simeq \cT M^k\,,$$
the diagonal prolongation $X^{[k]}$ of a vector field $X\in\X(M)$ can be understood as a vector field $ {X}^{[k]}$ on $M^k$, and the diagonal prolongation, $\alpha^{[k]}$, of a one-form $\alpha$ on $M$ can be understood as a one-form $ {\alpha}^{[k]}$ on $M^k$. If $k$ is assumed to be fixed, we will simply write $ {X}$ and $ {\alpha}$ for their diagonal prolongations. 

More explicitly,  if $Y$ be a vector field on a manifold $M$. The {\it diagonal prolongation} of $Y$ to $M^k$ is the vector field  
  \[ 
        {Y}^{[k]}(x_{(1)},\ldots,x_{(k)})=\sum_{a=1}^kY(x_{(a)})\]   on $M^k$ obtained by considering $\T M^k
       \simeq \T M\times \ldots \times \T M$ ($k$ times). Even more particularly, if $Y=x\partial/\partial x$ is a vector field on $\mathbb{R}$, then $Y^{[k]}=\sum_{a=1}^kx_{(a)}\partial/\partial x_{(a)}$.

The diagonal prolongation of a vector field can be extended to time-dependent vector fields on $M$, namely mappings $X:\mathbb{R}\times M\rightarrow \T M$ such that $X(t,\cdot)$ is a standard vector field $M$, by assuming that the diagonal extension to $M^k$ of the time-dependent vector field $X$ on $M$ is the time-dependent vector field $\widetilde{X}^{[k]}$ on $M^k$ whose value for every fixed valued of the time, $\widetilde{X}^k_t$, is the diagonal prolongation to $M^k$ of the vector field $X_t$. The space of diagonal prolongations of vector fields in $\mathfrak{X} (M)$ to $M^k$ form a Lie subalgebra of $\mathfrak{X}
(M^k)$. In fact, the mapping $X\in \X(M)\mapsto {X}^{[k]}\in \X(M^{k})$ is a Lie algebra morphism, i.e. it is a linear mappings such that
\begin{equation}
  \label{eq 14} [Y_1,Y_2]^{[k]}=\left[{Y}^{[k]}_1, {Y}^{[k]}_2\right],\qquad \forall Y_1,Y_2\in \mathfrak{X}(M)\,.
  \end{equation}

The following is a statement of the stochastic Lie theorem in \cite{LO09} without considering technical aspects that are not necessary for our purposes. It additionally solves a mistake in \cite{LO09}.

Let us solve a mistake in the proof of the stochastic Lie theorem in \cite{LO09}. The issue appears in the direct part of the statement \cite[p. 215, Theorem 3.1]{LO09}. For the sake of completeness, we will enunciate that part of the work adopting our notation and writing in full the contents referenced by the labels used in \cite{LO09}. We hereafter refer to
\begin{quote}\it
    ``Moreover, the fact
    $$
    \left[ {Z}^{[k+1]}_1,(Z_2+\lambda Z_3)^{[k+1]}\right]=([Z_1,Z_2+\lambda Z_3])^{[k+1]},\qquad \forall Z_1,Z_2\in \mathfrak{X}(M),\qquad \lambda\in \mathbb{R},
    $$
    and the hypothesis on $\{Y_1,\dotsc, Y_r\}$
    being in involution imply by the classical Frobenius Theorem that 
    $$
    D=\left\langle {Y}^{[k+1]}_1,\ldots,{Y}^{[k+1]}_r\right\rangle
    $$
    is integrable.''
\end{quote}
which is false. Let us explain the mistake and build a counterexample. Our counterexample will be relevant as it will show that certain stochastic differential equations in the It\^{o} approach are not stochastic Lie systems and do not accept a superposition rule. Note that $k$ is the smallest integer so that the vector fields $ {Y}^{[k]}_1,\ldots, {Y}^{[k]}_r$ are linearly independent at a generic point in $M^k$.

The problem relies on the fact that if $Y_1,\dotsc,Y_r$ span a distribution on $M$ that is involutive, then the diagonal prolongations $\widetilde Y^{[k+1]}_1,\ldots,\widetilde Y^{[k+1]}_r$ do not span the distribution given by the diagonal prolongations of the distribution spanned by $Y_1,\ldots,Y_r$. 

 
For instance, consider the two vector fields on $\mathbb{R}_\circ=\mathbb{R}\setminus\{0\}$ given by
$$
Y_1=x^2\frac{\partial}{\partial x}\,,\qquad Y_2=x^3\frac{\partial}{\partial x}
$$
that span an involutive distribution on $\mathbb{R}_\circ$. %(see \cite{Na66} for sufficient integrability conditions on distributions spanned by analytic vector fields). 
Notwithstanding, their diagonal prolongations, $ {Y}^{[s]}_1, {Y}^{[s]}_2$ to $\mathbb{R}_\circ^s$ , with $s>2$, do not need to span an involutive distribution. Indeed, one has the diagonal prolongations
$$
 {Y}_1^{[s]}=\sum_{a=1}^sx_{(a)}^2\frac{\partial}{\partial x_{(a)}},\qquad  {Y}_2^{[s]}=\sum_{a=1}^s x_{(a)}^3\frac{\partial}{\partial x_{(a)}}
$$
on $\mathbb{R}^s_\circ$. Their successive commutators become, as stated in \cite{LO09}, diagonal prolongations of an element of the involutive distribution spanned by $Y_1,Y_2$, namely $\T \mathbb{R}_\circ$. In particular,
$$
\ad_{ {Y}^{[s]}_1}^k {Y}^{[s]}_2=\left[ {Y}_1^{[s]},[
\ldots,[\ldots,[\ldots, [ {Y}_1^{[s]}, {Y}_2^{[s]}]\ldots]\ldots]\ldots]\right]=
\sum_{a=1}^sk!x_{(a)}^{3+k}\frac{\partial}{\partial x_{(a)}}=\left(k!x^{k+3}\frac{\partial}{\partial x}\right)^{[s]},
$$
for $k=1,2,\ldots$
Notwithstanding, any $\ad_{ {Y}^{[s]}_1}^k {Y}^{[s]}_2$, with $k\in \{2,3,4,\ldots\}$, is not contained in the distribution $\mathcal{D}=\left\langle  {Y}^{[s]}_1, {Y}^{[s]}_2\right\rangle$. Even worse, the smallest  (in the sense of inclusion) involutive  generalised distribution containing $ {Y}^{[s]}_1, {Y}^{[s]}_2$ spans the whole tangent space to $\mathbb{R}_\circ^s$ at almost every point and a superposition rule does not exist as it must be constructed from the non-constant first integrals of the vector fields taking values in an integral distribution containing  $ {Y}^{[s]}_1, {Y}^{[s]}_2$.  In fact, the vector fields $ {Y}^{[s]}_{k}=\sum_{a=1}^sx_{(a)}^{k+1}\frac{\partial}{\partial x_{(a)}}$, with $k=1,\ldots,s$, on $\mathbb{R}_\circ^s$ are linearly independent almost everywhere. To verify this fact, it is enough to see that the determinant of their coefficients read
$$
    \begin{vmatrix}
        x_{(1)}^2 & x_{(2)}^2 & \cdots & x_{(s)}^2 \\
        x_{(1)}^3 & x_{(2)}^3 & \cdots & x_{(s)}^3 \\
        \vdots & \vdots & & \vdots \\
        x_{(1)}^{s+1} & x_{(2)}^{s+1} & \cdots & x_{(s)}^{s+1}
    \end{vmatrix} = x_{(1)}^2\dotsb x_{(s)}^2 \begin{vmatrix}
        1 & 1 & \cdots & 1 \\
        x_{(1)} & x_{(2)} & \cdots & x_{(s)} \\
        \vdots & \vdots & & \vdots \\
        x_{(1)}^{s-1} & x_{(2)}^{s-1} & \cdots & x_{(s)}^{s-1}
    \end{vmatrix} = \prod_{a = 1}^s x_{(a)}^2 \prod_{1\leq i < j \leq s} (x_{(i)} - x_{(j)})\,.
$$
which makes that the smallest involutive generalised distribution containing $ {Y}^{[s]}_1, {Y}^{[s]}_2$ on $\mathbb{R}_\circ^s$ to be equal to $\T \mathbb{R}_\circ^s$ almost everywhere. This makes the existence of common non-constant first integrals for $ {Y}^{[s]}_1, {Y}^{[s]}_2$  impossible.

The previous counterexample is very important as it concerns the stochastic generalisations of the so-called Abel equations (see \cite{CLR11} and references therein). Hence, the mistake in \cite{LO09} may have potential consequences. In particular, note that the SIS model in the It\^{o} form \eqref{Eq:PseudoSto}, for a non-constant function $\sigma(t)$, takes the Stratonovich form \eqref{Eq:StrSIS}. If one tries to write the $t$-dependent coefficient for $\delta$ as a linear combination with $t$-dependent functions of a family of vector fields spanning a finite-dimensional Lie algebra, one finds that one has to obtain a finite-dimensional Lie algebra of vector fields containing $x^2\partial/\partial x,x^3\partial/\partial x$, which is impossible as shown above.

%Note also that \cite{LO09} has been cited many times so far.

Notwithstanding, for some reason the authors in \cite{LO09} state that local superposition rules are available if the distribution $D$ is not regular. As general differential equations are not generally concerned with this case, it is possible that the mistake may have had not many consequences.

\begin{theorem}
  [Stochastic Lie theorem]\label{teorema Lie}Let
  \begin{equation}
    \delta \Gamma = \mathfrak{S}(X, \Gamma) \delta X \label{eq1bis}
  \end{equation}
  be a stochastic differential equation on $M$, where $X :
  \mathbb{R}_+ \times \Omega \rightarrow \mathbb{R}^\ell$ is a given
  $\mathbb{R}^\ell$-valued semi-martingale and $\mathfrak{S}(x, z) : \T_x \mathbb{R}^\ell
  \to \T_z M$, for every $x
\in \mathbb{R}^\ell$ and $z\in M$, is a Stratonovich operator from
  $\R^\ell$ to $M$. Then, \eqref{eq1bis} admits a superposition rule if and only if
\begin{equation}\label{decomposition with the Ys2} \mathfrak{S}_j (X, \Gamma) = \sum_{\alpha = 1}^r b_j^\alpha (X)
      Y_\alpha (\Gamma)\,,\qquad j=1,\ldots,\ell\,,
    \end{equation}
    for every  $x\in \mathbb{R}^\ell$ and $\Gamma\in M$, where the vector fields $\{ Y_1, \ldots, Y_r \}$ span an $r$-dimensional Vessiot--Guldberg Lie algebra  on $M$.
\end{theorem}
\begin{proof} Let us prove the direct part of the theorem, as the converse is already proved in \cite{LO09}. One has that there always exists a number $m$ such that the diagonal prolongation of a basis $Y_1,\ldots,Y_r$ of the Vessiot--Guldberg Lie algebra of the Lie system reaches rank $r$ on a generic point. Let be $\mathcal{D}_0$ the distribution  spanned by vector fields $Y^{[m+1]}_1,\ldots,Y^{[m+1]}_r$ and the distribution has constant rank on an open subset of $M^{m+1}$. Since the Lie brackets $[Y_i,Y_j]$ are linear combinations of the vector fields of $V$, then their diagonal prolongation to $M^{m+1}$ is a liner combination of the $Y^{[m+1]}_1,\ldots,Y^{[m+1]}_r$ and $\mathcal{D}_0$ is integrable. Then, one can extend $\mathcal{D}_0$ to an integrable distribution $\mathcal{D}$ of corank $n$ in $M^{n+1}$, i.e. $\mathcal{D}_0\subset \mathcal{D}$ on every point of an open neighbourhood of a point in $M^{n+1}$. This can always be done in such a manner that the leaves of $\mathcal{D}$ project diffeomorphically onto an open subset of $M^n$ via the canonical projection $\pi:(\Gamma_{(1)},\ldots,\Gamma_{(m+1)})\in M^{m+1}\mapsto(\Gamma_{(1)},\ldots,\Gamma_{(m)})\in M^m$. The obtained expression is a local superposition rule for \eqref{eq1bis}.
\end{proof}


 

 

\section{Hamiltonian stochastic Lie systems}\label{Sec:Hamiltonian_Stochastic_Lie_Systems}

This section studies stochastic Lie systems admitting a Lie algebra of Hamiltonian vector fields relative to a geometric structure. To motivate our approach, let us consider a stochastic analogue of the SIS system studied in \cite{EFSZ22} given by
\begin{equation}\label{eq:SIS}
\begin{dcases}
 \dfrac{\d q}{\d t} = q\rho_0(t,B)-q^2-\frac{1}{p^2}\,,\\
 \dfrac{\d p}{\d t} = -p\rho_0(t,B)+2pq\,.
\end{dcases}
\end{equation}
Our stochastic generalisation is obtained by assuming that $\rho_0(t,z)$ depends, additionally, on a stochastic variable $B$. In this case,  the associated Stratonovich operator takes the form
$$
\mathfrak{S}(t,B,q,p)=\begin{pmatrix} q\rho_0(t,B)-q^2-\dfrac{1}{p^2}\\-p\rho_0(t,B)+2pq \end{pmatrix}\,,
$$
which allows us to write
$$
\delta\Gamma= \begin{pmatrix} \delta q\\\delta p \end{pmatrix} = \mathfrak{S}(t,B,q,p)\delta t\,.
$$
Moreover, one also has that $\mathfrak{S}(t,B,q,p)$ can be written as
$$
\mathfrak{S}(t,B,q,p)=\rho_0(t,B)Y_1+Y_2\,,
$$
where the vector fields $Y_1$ and $Y_2$ are given by
$$
Y_1=q\frac{\partial}{\partial q}-p\frac{\partial}{\partial p}\,,\qquad Y_2=-\left(q^2+\frac{1}{p^2}\right)\frac{\partial}{\partial q}+2qp\frac{\partial}{\partial p}\,.
$$
They close a finite-dimensional Lie algebra
$$
[Y_1,Y_2]=Y_2
$$
isomorphic to the non-Abelian two-dimensional Lie algebra $\mathfrak{h}_2$.

On the other hand, these vector fields are Hamiltonian relative to the standard symplectic structure
$\omega = \d q\wedge \d p$. In fact,
$$
\iota_{Y_1}\omega=qp\,,\qquad \iota_{Y_2}\omega = \d\left(-q^2p-\frac{1}{p}\right).
$$
Hence, we call \eqref{eq:SIS} a {\it Hamiltonian stochastic Lie system}.

Let us study now another possible stochastic SIS system studied in \cite{GGHMP11} of the form
\begin{align}
\delta S&=[-\beta SI+(\gamma +\mu)I]\delta t-\sigma SI\delta \mathcal{B}\,,\\
\delta I&=[\beta SI-(\mu+\gamma) I]\delta t+\sigma SI \delta \mathcal{B}\,,
\end{align}
where we consider a semi-martingale $X:\mathbb{R}_+\times \Omega\rightarrow \mathbb{R}^2$ and $B:\mathbb{R}_+\times \Omega\rightarrow \mathbb{R}$ is a Brownian motion. In this case, one can consider the vector fields on $\mathbb{R}^2$ given by
$$
\bar{X}_1= \frac{1}{N}X_1 = \frac{-SI}{S+I}\parder{}{S} + \frac{SI}{S+I}\parder{}{I}\,,\quad \bar{X}_2 = X_2 - \bar X_1 = \left(I - \frac{SI}{S+I}\right)\parder{}{S} - \left(I - \frac{SI}{S+I}\right)\parder{}{I}\,.
$$
Then,
$
[\bar X_1,\bar X_2]=\bar X_2\,.
$
Thus, the Stratonovich operator reads
\begin{equation}
    \mathfrak{S} = (\beta N\bar X_1 + (\gamma + \mu)(\bar X_2 - \bar X_1), \sigma N\bar X_1)\,.
\end{equation}
One can define a Lie algebra of Hamiltonian vector fields spanned by $\bar X_1,\bar X_2$ relative to the Jacobi structure on $\mathbb{R}^2$ given by $\Lambda=0$ and the Reeb vector field $R=\partial/\partial S-\partial/\partial I$. In this case, 
$$
\bar X_1=\frac{SI}{S+I}R\,,\qquad \bar X_2=\left(I-\frac{SI}{S+I}\right)R\,,
$$
which means that both vector fields are Hamiltonian ones.

In this case, one has that
$$
\mathfrak{S}_j(X,z)=\sum_{\alpha=1}^2 b_j^\alpha(z)Y_\alpha(z)\,,\qquad j=1,2,3\,.
$$
Finally, let us consider the stochastic Lotka--Volterra system
\begin{align}
\delta N_1 &= (b_1-a_1N_2)N_1\delta t+\sigma_1 N_1\delta \omega_1\,,\\
\delta N_2 &= (b_2-a_2N_1)N_2)\delta t+\sigma_2N_2\delta \omega_2\,,
\end{align}
where $\omega_1,\omega_2$ are Brownian motions
that one can find in \cite{Ar03}. This system can be studied through the vector fields 
$$
X_1=N_1\frac{\partial}{\partial N_1}+N_2\frac{\partial}{\partial N_2}\,,\qquad X_2=N_1N_2\frac{\partial}{\partial N_1}+N_1N_2\frac{\partial}{\partial N_2}\,,
$$
that close a two dimensional non-commutative Lie algebra. If you consider a symplectic form
$$
\omega=\frac{1}{N_1N_2}\d N_1\wedge \d N_2
$$
then the previous vector fields are Hamiltonian relative to it.

The latter appears as a stochastic Hamiltonian system that has been used in the study of SIR systems \cite{Ar03} and it has led to an analogue of the Lie systems in \cite{CFH23}. 

Note that the advantage of these systems is that we can model straightforwardly such systems without the deterministic approximations (like the ones in \cite{CFH23}).

Let us consider an harmonic oscillator on $\mathbb{R}^n$ with white noise $W_1$ (see \cite{HH11,Ha19} for analogues on $\mathbb{R}$) given by
\begin{equation}\label{eq:HHOWN}
\delta\begin{pmatrix}
    x_i \\ y_i
\end{pmatrix} = \begin{pmatrix}
    0 & \omega \\ -\omega & 0
\end{pmatrix} \begin{pmatrix}
    x_i \\ y_i
\end{pmatrix} \delta t + \begin{pmatrix}
    0 & 0 \\ 0 & -\sigma
\end{pmatrix} \begin{pmatrix}
    x_i \\ y_i
\end{pmatrix} \delta W_1\,,\qquad i=1,\ldots,n\,.
\end{equation}
 which is defined on $\cT\mathbb{R}^n$. That is, this is a special case of Example \ref{eq:WhiteOscil} by taking $k=0$. This much changes the properties of the system. It can be proved that \eqref{eq:HOWN} is not Hamiltonian relative to any known structure on $\mathbb{R}^2$ (cf. \cite{LS21}). But \eqref{eq:HHOWN} can now be related to a Vessiot--Guldberg Lie algebra given  on $\mathbb{R}^2$ given by
 $$
 X_1=\sum_{i=1}^n y_i\frac{\partial}{\partial x_i}-x_i\frac{\partial}{\partial y_i}
\,,\qquad X_2=\sum_{i=1}^n\frac{\partial}{\partial x_i}\,,\qquad X_{3}=\sum_{i=1}^n\frac{\partial}{\partial y_i}\,.
 $$
 Since,
 $$
 [X_1,X_2]=X_3\,,\qquad [X_1,X_3]=-X_2\,,\qquad [X_2,X_3]=0\,,
 $$
 this Lie algebra is the diagonal prolongation to $\mathbb{R}^5$ of the Vessiot--Guldberg Lie algebra $P_0\simeq \mathbb{R}\ltimes \mathbb{R}^2$  \cite{BBHLS15,LS21}, which is known to be a Lie algebra of Hamiltonian vector fields relative to the symplectic form
 $$
 \omega=\sum_{i=1}^n\d x_i\wedge \d y_i\,.
 $$
Hence, $X_1,X_2,X_3$ have Hamiltonian functions given by
 $$
 h_1 = \frac{1}{2}\sum_{i=1}^n(y_i^2+x_i^2)\,,\qquad h_2 = \sum_{i=1}^ny_i\,,\qquad h_3=-\sum_{i=1}^nx_i\,.
 $$
 relative to $\omega$. Note that $h_1,h_2,h_3$ also span, along with $h_4=1$, a Lie algebra of functions isomorphic that is a central extension of  the Heisenberg Lie algebra  $\mathfrak{h}_3$.
 
Now, \eqref{eq:HHOWN} can be related to  Stratonovich operator of the form
 $$
\mathfrak{S}(x ,y ,t,W)=\left(\omega X_1,-\sigma X_3\right),
 $$
 which turns \eqref{eq:HHOWN} into a stochastic Lie system and $V$ into an associated Vessiot--Guldberg Lie algebra P$_0$.


\begin{definition} A {\it Hamiltonian stochastic Lie system} on a manifold $M$ relative to a probability space $\Omega$ and stochastic variables $X:\mathbb{R}_+\times\Omega\rightarrow \mathbb{R}^\ell$ is a stochastic Lie system related to a series of Stratonovich operator $\mathfrak{H}(X,\Gamma)$ whose associated Vessiot--Guldberg Lie algebra $\langle Y_1,\ldots,Y_r\rangle$ consists of Hamiltonian vector fields relative to some geometric structure on $M$.
\end{definition}

Of course, the above definition means that there are Poisson Hamiltonian stochastic Lie systems, Jacobi stochastic Lie systems, contact stochastic Lie systems, etc.

\section{Stability for stochastic Hamiltonian systems}\label{Sec:SHS}
Let us recall the stability theory for stochastic differential equations and apply it to stochastic Hamiltonian systems. In particular, we will be specially interested in Hamiltonian stochastic Lie systems. For a survey on the theory for general stochastic differential equations and other related results see \cite{Ar74,LP08,Sa70}.

Consider a stochastic differential equation on $M$ of the particular form \begin{equation}
  \label{stochastic} \delta \Gamma =\mathfrak{S}_t (t,
  \Gamma)\delta t+ \mathfrak{S}_Y (t,
  \Gamma) \delta Y\,,
\end{equation}
where $X=(t,Y):\mathbb{R}_+ \times \Omega \rightarrow \R^\ell$ is an
$\R^\ell$-valued semi-martingale and 
$$\mathfrak{S} (t, \Gamma)=(\mathfrak{S}_t(t,\Gamma),\mathfrak{S}_Y(t,\Gamma)) : \T_X \mathbb{R}^\ell
\to \T_\Gamma M$$ is the associated Stratonovich operator. As in previous sections, it is assumed that initial conditions are values of $M$ which are taken with probability equal to one. Note that coefficients of the Stratonovich operator are assume to be functions depending only on the time and $M$. Stochastic differential equations of this type are common in the literature \cite{BCC_01}.

An {\it equilibrium point} for the stochastic differential equation \eqref{stochastic} is a point $\Gamma_e\in M$ such that
$$
\mathfrak{S} (t,
  \Gamma_e)=0\,,\qquad \forall t\in \mathbb{R}\,.
$$
Note that, in this case, the stochastic disturbance, which is described by $\mathfrak{S}_Y$, does not act at the equilibrium point $\Gamma_e0$. It can be stressed that a more general concept of equilibrium for Stratonovich differential equations with coefficients depending on stochastic processes can be found in \cite{LO08}.

% \begin{definition} {\bf (see \cite[Definition 11.2.6]{Ar03})}
%    An equilibrium point of \eqref{stochastic} is said to be {\it stochastically stable} if, for every $\epsilon>0$, one has
%    $$
%    \lim_{c\rightarrow 0}P\left[\sup_{t_0\leq t<\infty}|X_t(c)|\geq s\right]=0\,.
%    $$
%    Otherwise, it is said to be {\it stochastically unstable}.
% \end{definition}

\begin{definition}
Let $M$ be a manifold and let
\begin{equation}\label{eq:Stabil}
\delta \Gamma = \mathfrak{S}(t, \Gamma)\delta X
\end{equation}
be a Stratonovich stochastic differential equation with solutions $\Gamma:\mathbb{R}_+\times\Omega\rightarrow M$.
Given $\Gamma_0 \in  M$ and $s \in  \mathbb{R}$, denote by $\Gamma^{s,\Gamma_0}$ the particular solution of \eqref{eq:Stabil} such that $\Gamma_s^{s,\Gamma_0}(\omega)=\Gamma_0$ for all $\omega\in \Omega$. Let $\Gamma_e \in  M$ be an equilibrium of \eqref{eq:Stabil}, that is, the constant process
$\Gamma_t(\omega) = \Gamma_e$, for all $t \in \mathbb{R}$ and $\omega\in \Omega$, is a solution of \eqref{eq:Stabil}. The equilibrium point $\Gamma_e$ is:
\begin{itemize}
\item Almost surely (Lyapunov) stable if for  any open neighbourhood $U$ of $\Gamma_e$ there exists a
neighbourhood $V$ of $\Gamma_e$ such that for any $\Gamma'\in V$, one has that $\Gamma^{s,\Gamma'}\subset U$ almost surely (a.s.)
\item Stable in probability. For any $s\geq 0$ and $\rho>0$ one has that
$$
\lim_{x\rightarrow {z_0}} P\left\{\sup_{t>s}d(\Gamma^{s,\Gamma_0}_t,\Gamma_0)>\epsilon\right\}=0,
$$
where $d:M\times M\to\mathbb{R}$ is any distance function that generates the manifold topology.
\end{itemize}
\end{definition}

Let us provide a relevant analogue of the method used in Hamiltonian symplectic systems to check stability \cite{AM78}.
%\begin{definition}
 % Let $U$ be an open neighbourhood
%of the equilibrium $z_0$ and let $V\in \Cinfty(U)$. Then, $ V$ is a Lyapunov
%function for the equilibrium $z_0$ if $V (z_0) = 0$, and  $V (z) > 0$ for any $z \in U \setminus\{z_0\}$ , and
%$$
%E[V (\Gamma_\tau )] \leq  E[V (\Gamma_0)]
%$$
%for any stopping time $\tau$ and any solution $\Gamma$ of \eqref{eq:Stabil}.    
%\end{definition}
The most relevant result for our purposes in the study of stability of Hamiltonian stochastic Lie systems is the following proposition. 
\begin{proposition} (Stochastic Dirichlet’s Criterion) Suppose that we are in the setup of the previous definition and assume that there exists a function $f\in \Cinfty(M)$ such that $\d f_{z_0} = 0$ and that the
quadratic form $\d^2f_{z_0}$  (positive or negative) definite. If $f$ is a strongly (respectively, weakly) conserved quantity for the solutions of \eqref{eq:Stabil} then the equilibrium $z_0$ is almost surely stable (respectively,
stable in probability.
\end{proposition}

Let us turn now to Hamiltonian stochastic differential equations on symplectic manifolds. In this case, we focus on stochastic differential equations of the type
\begin{equation}\label{eq:HamSDE}
\delta  \Gamma=\mathfrak{H}(X,\Gamma )\delta X\,,
\end{equation}
where $X:\mathbb{R}_+\times \Omega\rightarrow \mathbb{R}^\ell$ is a semi-martingale and $\mathfrak{H}:\T \mathbb{R}^\ell\times M\rightarrow \T M$ is a Stratonovich operator such that
$$
\mathfrak{H}(X,\Gamma)=(X_1(\Gamma),\ldots,X_\ell(\Gamma))\,, 
$$
where $X_1,\ldots,X_\ell$ stand for Hamiltonian vector fields relative to a symplectic form $\omega$ on $M$ with Hamiltonian functions $h_1,\ldots,h_r\in \Cinfty(M)$, respectively. 

The symplectic structure in this case allows for the study of the system via the Hamiltonian functions of the system, which gives powerful methods to study their properties. Note that \eqref{eq:HamSDE} has indeed associated a $h:\mathbb{R}_+\times \Omega\rightarrow \mathbb{R}^{\ell*}$ of the form
$$
h=(h_1,\ldots,h_\ell).
$$
Normally, we call $h_1$ the Hamiltonian of the system, but Hamiltonian stochastic differential equations have several ones. In fact, note that $h_1$ need not be conserved.


It is worth noting that we are mainly interested in Hamiltonian stochastic Lie systems, e.g.
\begin{equation}\label{eq:HamStoLieSys}
\delta \Gamma=\sum_{\alpha=1}^rb_\alpha(t,  Y)X_\alpha \delta t+\sum_{i=2}^\ell\sum_{\alpha=1}^rb^i_\alpha(t, Y)X_\alpha\delta Y^i
\end{equation}
such that the vector fields $X_1,\ldots,X_r$ span an $r$-dimensional Vessiot--Guldberg Lie algebra of Hamiltonian vector fields relative to a geometric structure structure. In particular, we here focus on Hamiltonian Lie systems relative to a symplectic structure whose coefficients in \eqref{eq:HamStoLieSys} depend only on time. In such a case, \eqref{eq:HamStoLieSys} is then related to a series of $\ell$-Hamiltonian functions being a linear combination with coefficients depending on time of certain Hamiltonian functions $h_1,\ldots,h_r$ spanning a Lie algebra of Hamiltonian functions relative to the symplectic structure. In particular, one will have
$$
h_0=\sum_{\alpha=1}^rb^\alpha_0(t)h_\alpha\,,\qquad h_i=\sum_{\alpha=1}^rb_\alpha^i(t)h_\alpha.
$$
It is standard to call $h_0$ the Hamiltonian of the stochastic differential equation. It is immediate that it is not a constant of the motion in general as $\{h_0,h_i\}$ not need to be zero in general. 

Notwithstanding, it is remarkable that a function on $M$ will be a constant of motion for \eqref{eq:HamStoLieSys} if 
$$
\{h_\alpha,f\}=0\,,\qquad \alpha=1,\ldots,r\,.
$$
Let us assume that there exists a function $f$ that is a constant of motion of $\delta \Gamma$ and it also satisfies that it has a minimum at an equilibrium point of $\delta \Gamma$. Then, it is immediate that $\delta\Gamma$ has a stable equilibrium point.

It is simple to apply the above results to linear or affine stochastic differential equations, which are stochastic Lie systems, and to choose cases with a Vessiot--Guldberg Lie algebra of Hamiltonian vector fields relative to a symplectic structure.

\section{Relative equilibrium points for stochastic Hamiltonian Lie systems}\label{Sec:Relative_Equilibrium_Stochastic_Differential}

In short, relative equilibrium points for a differential equation, in general, and stochastic one, in particular, are points where the dynamics is generated by symmetries of the differential equations. In this section we will define relative equilibrium points for Hamiltonian stochastic differential equations in a Stratonovich form. This will be developed in terms of the theory developed in \cite{LO08}. 

Let us start be reviewing the notion of symmetries and symplectic reduction for Hamiltonian stochastic differential equations. 

\begin{definition}
Let $X:\mathbb{R}_+\times \Omega\rightarrow M$ be a semi-martingale and let $\mathfrak{S}:\T\mathbb{R}^\ell\times M\rightarrow \T M$ be a Stratonovich operator. A diffeomorphism $\phi:M\rightarrow M$ is a symmetry of the stochastic differential equation associated with $\mathfrak{S}$ if
$$
\mathfrak{S}(v,\phi (\Gamma))=\T\phi [\mathfrak{S}(v,\Gamma)],\qquad \forall (v,\Gamma)\in \T \mathbb{R}^\ell\times M.
$$
More generally, a Lie group action $\Phi:G\times M\rightarrow M$ is a symmetry of  the stochastic differential equation induced by $\mathfrak{S}$ if $\Phi_g$ is a symmetry of $\mathfrak{S}$ for every $g\in G$. Similarly, we say that $\phi$ and $\Phi$ are symmetries of the Stratonovich operator $\mathfrak{S}$.     
\end{definition}

The relevance of symmetries of Stratonovich operators is due to the fact that they transform a particular solution of the stochastic differential equations related to the Stratonovich operator to another particular solution of the same stochastic differential equation.

Let us review the stochastic Marsden--Weinstein reduction for Stochastic Hamiltonian systems as introduced in  \cite{LO9Red}. 
 Let  $(M,\omega)$ be a symplectic manifold and let $\Phi:G\times M\rightarrow M$ be a Lie group action  admitting a coadjoint equivariant momentum map $J : M \rightarrow \mathfrak{g}^*$ being a Lie group symmetry of the Hamiltonian functions $h:M\rightarrow \mathbb{R}^\ell$ of a Stratonovich operator $\mathfrak{H}
$. Every  regular $\mu\in \mathfrak{g}^*$ of $J$ gives rise to a   function  $h_\mu  :  J^{-1}(\mu)/G_\mu \rightarrow \mathbb{R}^\ell$ on the manifold $M_\mu=J^{-1}(\mu)/G_\mu$
determined by the equality $h_\mu\circ \pi_\mu=h\circ \iota_\mu$, where $\iota_\mu:J^{-1}(\mu)\hookrightarrow M$ is the natural immersion. In turn, this induces a stochastic Hamiltonian
system on the symplectic reduced space $(M_\mu := J^{-1}(\mu)/G_\mu, \omega_\mu)$ with stochastic
component $X$ and whose Stratonovich operator $\mathfrak{H}_\mu : \T \mathbb{R}^\ell \times M_\mu \to \T M_\mu$ is
given by
$$
\mathfrak{H}_\mu(t, [\Gamma]))(u) = \left(\sum_{\alpha=1}^rb_\alpha(t)X_{h^\mu_\alpha}([\Gamma]),\sum_{\alpha=1}^rb_\alpha(t)X_{h^\mu_\alpha}([\Gamma])\right),
$$
where
$h^\mu_i \in\Cinfty(M_\mu)$ are the coefficients of $h^\mu:[\Gamma]\in M_\mu\mapsto (h^\mu_1(\Gamma),\ldots,h^\mu_r(\Gamma))$ for every $\Gamma\in J^{-1}(\mu)$ induced  via the relation $h_\mu \circ \pi_\mu = h \circ \iota_\mu$. Moreover, if $\Gamma$ is a solution semi-martingale of the Hamiltonian
system associated with $\mathfrak{H}$ with initial condition $\Gamma_0\subset J^{-1}(\mu)$, then so is $\Gamma_\mu := \pi_\mu(\Gamma)$
with respect to $\mathfrak{H}_\mu$, with initial condition $\pi_\mu(\Gamma_0)$.

Then, one has the following definition.

\begin{definition} Given a Hamiltonian  stochastic differential equation \eqref{eq:HamSDE} on $M$, a {\it relative equilibrium point}  $\Gamma_r\in M$ of \eqref{eq:HamSDE} is a point such that
\begin{equation}\label{eq:RelEquH}
\mathfrak{H}(t,\Gamma_e)\in D_x,\qquad  
\end{equation}
where $D$ is the distribution generated by the fundamental vector fields of a Lie group action $\Phi:G\times M\rightarrow M$ with a momentum map $J:M\rightarrow \mathfrak{g}^*$.
\end{definition}

To understand the above condition in terms of the better known characterisation for standard Hamiltonian systems of differential equations \cite{MS88}, let   $\xi_M^1,\ldots, \xi^r_M$ be a basis of fundamental vector fields of $\Phi$ corresponding to a basis  $\xi_1,\ldots,\xi_r$ of $\mathfrak{g}$. Then, the relation \eqref{eq:RelEquH} amounts to saying that
$$
\mathfrak{H}(e,\Gamma)=\sum_{\alpha=1}^rf^\alpha (e)\xi_\alpha^M(\Gamma)\,.
$$
for certain functions $f^1,\ldots,f^r\in \Cinfty(V)$. 

Let $\mu=J(x)$ in the above definition. If the Hamiltonian stochastic differential equation \eqref{eq:HamSDE} can be reduced, then $\Gamma\in M$ is contained in $J^{-1}(\mu)$. If it is assumed to be a submanifold of $M$. The evolution of $H$ is contained within the leaves of $J^{-1}(\mu)$. This implies that the right-hand side \eqref{eq:RelEquH} belongs indeed to $D_x\cap \T_xJ^{-1}(\mu)$ Moreover, the restriction of $H$ to it can be reduced to  $J^{-1}(\mu)/G_\mu$. Moreover, the $\mathfrak{H}_\mu(e,\pi_\mu)=0$. In other words, one has the following proposition, whose inverse is analogue to its deterministic case.

\begin{proposition} Let $\Gamma\in M$ be a relative equilibrium point for a stochastic Hamiltonian system \eqref{eq:HamSDE} and let $J(x)=\mu$ be a regular point of $J$ so that $J^{-1}(\mu)/G_\mu$ is a manifold. Then, the reduction of \eqref{eq:RelEquH} to $M_\mu=J^{-1}(\mu)/G_\mu$  is such that $\pi_\mu(\Gamma)$ is an equilibrium point of the reduced system and vice versa.
\end{proposition}

Finally, let us characterise relative equilibrium points for stochastic Hamiltonian systems.

\begin{theorem}\label{RelativeEquilibrum}
{\bf (Stochastic Relative Equilibrium Theorem)} A point $\Gamma_{e}\in M$ is a relative equilibrium point for a Hamiltonian stochastic differential equations \eqref{eq:HamStoLieSys} if and only if there exist elements $\xi_e=(\xi^1,\ldots,\xi^\ell)\in\mathfrak{g}^\ell$ such that $(\xi_e,\Gamma_{e})$ is a critical point of the coordinates of the function $\widehat{h}:\mathfrak{g}^\ell\times M\rightarrow \mathbb{R}^\ell$ given by
\[\widehat{h}^\alpha(\Gamma,\xi)=h^\alpha(\Gamma)-\left\langle {J}(\Gamma) - \mu_{e},\ \xi^\alpha\right\rangle,\qquad \alpha=1,\ldots,\ell,
\] 
for $\mu_e:={\bf J}(\Gamma_e)$.
\end{theorem}
\begin{proof} At the critical point $(\xi_e,\Gamma_e)$,  the function $\widehat{\Gamma}(\xi_e,\cdot):M\rightarrow \mathbb{R}^\ell$ has a critical point $\Gamma_e\in M$. Hence, each coordinate of $\widehat{h}(\xi_e,\cdot)$ satisfies that $d(h^\alpha-\langle J(\cdot),\xi^\alpha_e\rangle)|_{\Gamma_e}=0$. This implies that $X_{h^\alpha}|_{\Gamma_e}=X_{J_{\xi^\alpha_e}}|_{\Gamma_e}$ and the point $\Gamma_e$ is a relative equilibrium point of $\mathfrak{H}$.  Conversely, $\Gamma_e$ is a relative equilibrium point to $\mathfrak{H}$, at each coordinate of $\mathfrak{H}$, one has that $X_{h_\alpha}=\xi_M^\alpha$ for certain $\xi^1,\ldots,\xi^\ell\in \mathfrak{g}$. Consequently, one has that $X_{h^\alpha-J^\alpha}=0$. This implies that each one of these vector fields has a Hamiltonian function given by $h^\alpha-\langle J-\mu_e,\xi^\alpha\rangle$ whose has a critical point at the given point. It is worth noting that $(\xi_e,\Gamma_e)$ is a equilibrium point of $\widehat{h}$.  
\end{proof}



\section{Stochastic Poisson coalgebra method}\label{Se:PCSLS}

Let us review the Poisson coalgebra method for Hamiltonian systems relative to different geometric structures. In general, our approach slightly matches what can be found in previous works, but it also introduces certain simplifications and assumptions that allow one to apply it to many topics. In particular, the method can still be applied to Hamiltonian stochastic Lie systems by considering that they are determined by a certain family of Hamiltonian functions.


Let us provide a method to derive superposition rules for general Hamiltonian stochastic Lie systems via Poisson coalgebras. Our present procedure is a new modification of the coalgebra method for deriving superposition rules for $k$-symplectic Lie systems devised in \cite{LV15}. 

It is convenient to stress that the proof of the Stochastic Lie theorem shows that a superposition rule for stochastic Lie systems, in Stratonovich form, can be obtained in a similar manner to the case of deterministic Lie systems, namely
	
	The procedure to obtain a superposition rule goes as follows (see \cite{CGM07,Dissertationes,LS21} for
	details and examples).
	\begin{enumerate}
		\item Consider a Vessiot--Guldberg Lie algebra spanned by a basis of vector fields $X_1,\ldots,X_r$ on the manifold $M$ where the stochastic Lie system is defined on.
  \item Find the smallest  $m\in \mathbb{M}$, so that $X^{[m]}_1,\ldots,X_r^{[m]}$  
		are linearly independent at a generic point.
		\item Use local coordinates $x^1,\ldots,x^n$ on $M$ and consider this coordinate
		system to be defined on each copy of $M$ within $M^{m+1}$ to get a coordinate system
		$\{x^i_{(a)}\mid i=1,\ldots,n,\,\,a=0,\ldots,m\}$ on $M^{m+1}$.
		Obtain $n$ functionally independent first-integrals $F_1,\ldots, F_n$ common to all diagonal prolongations
		$X^{[m+1]}_1,\ldots,  X^{[m+1]}_r$ such that 
		\begin{equation}\label{Cond}
		\frac{\partial(F_1,\ldots,F_n)}{\partial(x_{(0)}^1,\ldots,x_{(0)}^n)}\neq 0.
		\end{equation}
		  
		\item 
	    The condition \eqref{Cond} allows us to ensure that the equations $F_i=k_i$, for $i=1,\ldots,n$, enable us to write the expressions of the variables $x_{(0)}^1\ldots,x_{(0)}^n$ in
		terms of $x^1_{(a)},\ldots,x_{(a)}^n$, with $a=1,\ldots,m,$  and $k_1,\ldots,k_n$.
		\item The obtained expressions lead to a superposition rule depending on a generic family of $m$ particular solutions
		and the constants $k_1,\ldots, k_n$.  
	\end{enumerate}


Note that every Hamiltonian stochastic Lie system is related to a Stratonovich operator $\mathfrak{H}$ given by $\ell$ components and each one can be understood as a vector field. Hence, $\mathfrak{H}$ can be understood as a $\ell$-vector field, i.e. a section of the $\ell$-tangent bundle of velocities $\T^\ell M=\T M\oplus \dotsb \oplus \T M\rightarrow M$, where $\T M\oplus\dotsb\oplus \T M$ is to be considered as the Whitney sum of the tangent bundle to $M$ with itself. In this manner, one can prolong it to a section $\mathfrak{H}$ of $(\T^kM)^{[m]}\rightarrow M^m$. 
Moreover, $\mathfrak{H}$ is related to a certain Hamiltonian function $h:M\rightarrow \mathbb{R}^\ell$. Then, $h$ is a section of $M\times \mathbb{R}^\ell M$ and it can be extended to a section of $(M\times \mathbb{R}^\ell)^m\rightarrow M^m$. Moreover, the prolongation $h^{[m]}$ is the Hamiltonian function related to $\mathfrak{H}^{[m]}$.  Then, the following result is immediate.

\begin{proposition}\label{Cor:LieAlgJL2}
    If $\mathfrak{H}$ is a Hamiltonian stochastic Lie system admitting a Hamiltonian $h:M\rightarrow \mathbb{R}^\ell$ relative to $\omega$, then $\mathfrak{H}^{[m]}$ is a Hamiltonian stochastic Lie system relative to $\omega^{[m]}$ admitting a Hamiltonian $h^{[m]}:M^m\rightarrow \mathbb{R}^\ell$.
\end{proposition}

One has the following immediate result.
\begin{proposition} The space of $\ell$-Hamiltonian functions on $M$ is a Poisson algebra relative to the bracket
$$
\{h,h'\}_\ell=(\{h_1,h'_1\}_\omega,\ldots,\{h_\ell,h'_\ell\}_\omega)\,, \qquad \forall h,h'\in \Cinfty(M,\mathbb{R}^\ell)\,,
$$
and the multiplication
$$
h\cdot h'=(h_1h_1',\ldots,h_\ell h'_\ell)\,,\qquad 
\forall h,h'\in \Cinfty(M,\mathbb{R}^\ell)\,.
$$
\end{proposition}
It is worth recalling that the case of $\ell$-Hamiltonian functions relative to an $\ell$-symplectic structure was not a Poisson algebra due to the fact that the multiplication of $\ell$-Hamiltonian functions could not be ensured to be an $\ell$-Hamiltonian function as each coordinate is assumed to be a Hamiltonian function of the same vector field and this cannot be ensured as each coordinate of the $\ell$-function is associated with different presymplectic structures.

%In other words, $[R,\Lambda]_{SN}=0$ and $[\Lambda,\Lambda]_{SN}=2R\wedge \Lambda$. Let us describe a coalgebra method for Jacobi--Lie systems \cite{AHR22,AHKR22,LS2020}. %In turn, this gives rise to a subbundle of the Pontryagin bundle $\mathbb{T}N=TN\oplus_NT^*N$ spanned by \cite{Cou1990}
%$$
%\iota_{du^1}\Lambda+du^1,\ldots,\iota_{du^{n-1}}\Lambda+du^{n-1},R
%$$
%where $u^1,\ldots,u^r$ are functionally independent first integrals of $R$. 

%The proofs of Proposition \ref{Pro:ProlJL}, its Corollaries \ref{Cor:LieAlgJL} and \ref{Cor:LieAlgJL2}, and Proposition \ref{Prop:Obv2} below are straightforward as they rely, almost entirely, on the definition of diagonal prolongations. Anyhow, as Jacobi manifolds with a non-vanishing Reeb vector field give rise to a Dirac manifold, they can also be considered as slight modifications of the results given for Dirac structures in \cite{Car2014}. %and the fact that the diagonal prolongation to $N^m$ of the Schouten--Nijenhuis bracket of two multivector fields on $N$ is equal to the Nijenhuis bracket of the diagonal prolongations to $N^m$ of the vector fields on $N$.  %Moreover, one has to recall that every function $f\in \Cinfty(N)$ admits an associated Hamiltonian vector field relative to $(N,\Lambda,R)$ given by $\Lambda(df)+fR$.


\begin{proposition} Let $\omega$ be a symplectic form on $M$ and let $\{e^1,\ldots,e^k\}$ be a basis of $\mathbb{R}^\ell$.  Then, the mapping $\phi_\theta:(h_1,\ldots,h_\ell)\in \Cinfty(M)\otimes \mathbb{R}^\ell\mapsto \sum_{i=1}^\ell\theta_ih_i\in \Cinfty(M)$ is a Lie algebra morphism if $\Cinfty(M)$ is endowed with bracket of $\omega$.
\end{proposition}

% \begin{proposition}\label{Pro:ProlJL} Let $(M,\Lambda,E)$ be a Jacobi manifold with bracket $\{\cdot,\cdot\}_{\Lambda,E}$. Let $X$ and $f$ be a vector field and a function on $M$, respectively. Then,
% \begin{enumerate}[(a)]
%     \item $(M^k,\Lambda^{[k]},E^{[k]})$ is a Jacobi manifold for every $k\in \mathbb{N}$.
%     \item If $f$ is a Hamiltonian function for a Hamiltonian vector field $X$ relative to $(M,\Lambda,E)$, its diagonal prolongation $f^{[k]}$ to $M^k$ is a Hamiltonian function of the diagonal prolongation, $X^{[k]}$, to $M^k$ with respect to $(M^k,\Lambda^{[k]},E^{[k]})$. 
%     \item If $f \in {\rm Cas}(M ,\Lambda, E)$, then $f^{[k]}\in {\rm Cas}(M^k,\Lambda^{[k]}, E^{[k]})$.
%     \item The map $\lambda: (\Cinfty(M),\{\cdot,\cdot\}_{\Lambda,E})\ni f \mapsto f^{[k]}\in (\Cinfty(M^k), \{\cdot,\cdot \}_{\Lambda^{[k]},E^{[k]}})$ is an injective Lie algebra morphism.
% \end{enumerate}
% \end{proposition}
 
%\begin{corollary}\label{Cor:LieAlgJL}
 %   Let $h_1,\dotsc, h_r:M\to \mathbb{R}$ be a family of functions on a Jacobi manifold $(M,\Lambda,E)$ spanning a finite-dimensional real Lie algebra of functions with respect to the Lie bracket $\{\cdot,\cdot\}_{\Lambda,E}$. Then, their diagonal prolongations to $M^k$, $\tilde h_1, \dotsc,\tilde h_r$, close an isomorphic Lie algebra of functions with respect to the Lie bracket $\{\cdot,\cdot\}_{\Lambda,E}$ induced by the Jacobi manifold $(M^k,\Lambda^{[k]}, E^{[k]})$.
%\end{corollary}



%\begin{proposition}\label{Prop:Obv2}
 %   If $X$ be a system possessing a time-independent constant of the motion $f$ and $Y$ is a time-independent Lie symmetry of $X$, then:
  %  \begin{enumerate}
   %     \item The diagonal prolongation $f^{[k]}$ is a time-independent constant of the motion for $X^{[k]}$.
    %    \item The vector field $Y^{[k]}$ is a time-independent Lie symmetry of $X^{[k]}$.
     %   \item If $h$ is a time-independent constant of the motion for $X^{[k]}$, then $Y^{[k]}h$ is another time-independent constant of the motion for $X^{[k]}$.
    %\end{enumerate}
%\end{proposition}
%\begin{proposition} The diagonal prolongation to $M^k$ of a Jacobi--Lie system $(M,\Lambda,E,X)$ is a Jacobi--Lie system $(M^k,\Lambda^{[k]},E^{[k]},X^{[k]})$.
%\end{proposition}

Note that the diagonal prolongation of a contact Lie system do not need to be a contact Lie system as it may not be defined on an odd-dimensional manifold.

For the sake of completeness, let us prove the following result.

\begin{proposition}
    Let $\mathfrak{H}$ be a Hamiltonian stochastic Lie system with respect to a symplectic form $\omega$ and possessing a smallest (in the sense of inclusion) Lie--Hamiltonian Vessiot--Guldberg Lie algebra $(\mathfrak{W},\{h_1,\ldots,h_r\})$ relative to the Poisson bracket related to $\omega$. A function $f\in \Cinfty(M)$ is a constant of the motion for $\mathfrak{H}$ if and only if it commutes with all the elements of ${\rm Lie}(\{h_t\}_{t\in \R},\{\cdot,\cdot\})$.
\end{proposition}
\begin{proof} The function $f$ is a constant of the motion for $X$ if
\begin{equation}\label{Eq:Cons}
0=X_tf=\{h_t,f\}\,,\qquad \forall t\in \mathbb{R}\,.
\end{equation}
Hence, 
$$
\{f,\{h_t,h_{t'}\}\}=\{\{f,h_t\},h_{t'}\}\}+\{h_t,\{f,h_{t'}\}\}\,,\qquad \forall t,t'\in \mathbb{R}\,.
$$
Inductively, $f$ is shown to commute with all the elements of ${\rm Lie}(\{h_t\}_{t\in \mathbb{R}})$. 

Conversely, if $f$ commutes with all ${\rm Lie}(\{h_t\}_{t\in \mathbb{R}})$ relative to $\{\cdot,\cdot\}$, in particular, \eqref{Eq:Cons} holds and $f$ is a constant of the motion of $X$.
\end{proof}

In Lie--Hamilton systems, a $t$-dependent Hamiltonian vector field admitted for every $t\in \mathbb{R}$ a Hamiltonian function belonging to a finite-dimensional Lie algebra of Hamiltonian functions relative to a Poisson bracket. 

%The bracket for Jacobi manifolds is not, in general, a Poisson bracket. It becomes only a Poisson bracket for good Hamiltonian functions. Nevertheless, when a Lie group action gives rise to a momentum map, the components of the momentum map are first integrals of $R$. As a consequence, the following proposition, which can be considered as an adaptation of \cite[Proposition 8.4]{Car2014}, is satisfied. Recall that if $(M,\Lambda,E)$ is a Jacobi manifold, then $\Cinfty(M^k)$ becomes a Lie algebra relative to the Lie bracket $\{\cdot,\cdot\}_k$ related to $\Lambda^{[k]}$ and $\Cinfty(\mathfrak{W}^*)$ is a Poisson algebra relative to the Kirillov--Kostant--Souriau bracket.

\begin{proposition}\label{Prop:6.7}
    Let $\mathfrak{H}$ be a Hamiltonian stochastic Lie system with an associated Lie--Hamilton Vessiot--Guldberg Lie algebra $(\mathfrak{W}=\langle h_1,\ldots,h_r\rangle, \{\cdot,\cdot\}_\omega)$ relative to the symplectic form $\omega$. Let $\{v_1,\ldots v_r\}$ be a basis of linear coordinates on $\mathfrak{W}^*$. Given the good momentum map $J : M\to \mathfrak{W}^\ast$, the pull-back $J^\ast C$ of any Casimir function $C$ on $\mathfrak{W}^\ast$ is a constant of the motion for $\mathfrak{H}$. Moreover, if  $h_i=J^*v_i$ for $i=1,\ldots,r$, and $C=C(v_1,\ldots,v_r)$, then 
    \begin{equation}\label{Eq:ExCas}
C\left(\sum_{a=1}^kh_1(x_{(a)}),\ldots, \sum_{a=1}^kh_r(x_{(a)})\right)\,,
\end{equation} is a constant of the motion of $\mathfrak{H}^{[k]}$.
\end{proposition}

The stochastic Poisson coalgebra method takes its name from the fact that it is applied to Hamiltonian stochastic Lie systems and it analyses the use of Poisson coalgebras and a so-called {\it coproduct} to obtain superposition rules. In fact, the coproduct is responsible for the form of \eqref{Eq:ExCas}.


% in this theory is a map $\Delta:f(v_1,\ldots,v_r)\in \Cinfty(\mathfrak{g}^*)\mapsto f(v_1^{(1)}+v_1^{(2)},\ldots,v_r^{(1)}+v_r^{(2)})\in  \Cinfty(\mathfrak{g}^*)\otimes \Cinfty(\mathfrak{g}^*)$ that is co-commutative, i.e. ${\rm Id}\otimes \Delta=\Delta \otimes {
%\rm Id}$. In turn, one can determine a mapping $\Delta^{k}:\Cinfty(\mathfrak{g}^*)\rightarrow \Cinfty(\mathfrak{g}^*)\otimes\ldots\otimes\Cinfty(\mathfrak{g}^*)$ ($k$-times).  Moreover, every contact Lie system gives rise to a mapping $D:f\in \mathfrak{W}\mapsto X_f\in \mathfrak{X}_{ham}(M)$. It can be proved that if $C$ is a Casimir of $\mathfrak{g}$, then $D^m\Delta(C)$ is a Casimir function on $\mathfrak{W}^*$, which can be used to obtain a superposition rule. 

\section{Conclusions and outlook}\label{Eq:Conclusions}

This paper has reviewed in a simple manner the theory of stochastic Lie systems solving a mistake in the literature. Some relations between the Stratonovich and the It\^o approaches to stochastic Lie systems have been reviewed, which has shown that the stochastic Lie theorem in the It\^o approach has a different non-expected form. Types of Hamiltonian stochastic Lie systems have been introduced as well as some possible generalisations. We have also introduced the theory of Hamilton stochastic Lie systems. The theory of stability and relative equilibrium points for Hamiltonian stochastic differential equations has been developed and, in particular, we have focused in the study of Hamiltonian stochastic Lie systems relative to symplectic forms. Some results concerning the energy-momentum method for stochastic Lie systems have been obtained, generalizing \cite{Bz_14}. Several examples concerning physical applications have been developed. We have also remarked that the theory of Poisson coalgebras can be developed for our Hamiltonian stochastic Lie systems.

Concerning applications, several stochastic SIS models have been analysed. Such models use to be studied using deterministic methods in the literature. Instead, we here approach them in a pure stochastic approach. 

In the future, we aim to further develop the energy-momentum method for stochastic Lie systems. This entails the search for criteria ensuring our of a Hamiltonian Stratonovich operator the stability of equilibrium points obtained after reduction. Moreover, we expect, as in the deterministic case, to obtain certain degeneracies that will have to be analysed to solve the problem.  Note that a further study of the advantages of epidemiological models from a pure stochastic point of view will have to be addressed. Moreover, we also plan to study a stochastic Hamiltonian system appearing in celestial mechanics concerning the stochastic variation of the inertia tensor \cite{Sa70}.



\section*{Acknowledgements}
\addcontentsline{toc}{section}{Acknowledgements}
Fruitful discussions with J.P. Ortega to fix a correction of his stochastic Lie theorem are acknowledged. 
%X. Rivas acknowledge  financial support from the Nowe Idee 2B-POB II project PSP: 501-D111-20-2004310 funded by the ``Inicjatywa Doskonałości - Uczelnia Badawcza'' (IDUB) program. 
X. Rivas acknowledges partial financial support from the Spanish Ministry of Science and Innovation grants PID2021-125515NB-C21 and RED2022-134301-T of AEI, and the Ministry of Research and Universities of the Catalan Government project 2021 SGR 00603.

J. de Lucas acknowledges a MINIATURA-5 grant by the Polish Narodowe Centrum Nauki under project Nr 2021/05/X/ST1/01797. %Fruitful conversations with J.P. Ortega are acknowledged.


The research of Marcin Zaj\k ac was funded by the Polish National Science Centre under the contract number UMO-2021/41/N/ST1/02908. For the purpose of Open Access, the author has applied a CC-BY public copyright licence to any
Author Accepted Manuscript (AAM) version arising from this submission.

%{\addcontentsline{toc}{section}{Bibliography}}

	
\bibliographystyle{abbrv}
%%%%%%%%%%%%%%%%%%%%%%%%%%%%%%%%%%%%%%%%%%%%%%%%%%%%%%%%%%%%%%%%
%%%%%%%%%%%%%%%%%%%%%%%%%%%%%%%%%%%%%%%%%%%%%%%%%%%%%%%%%%%%%%%%
\addcontentsline{toc}{section}{References}
\itemsep 0pt plus 1pt
%\small
\bibliography{references.bib}

\end{document}


\begin{thebibliography}{10}

\bibitem{AM78}
R.~Abraham and J.~E. Marsden.
\newblock {\em Foundations of mechanics}.
\newblock Benjamin/Cummings Publishing Co., Inc., Advanced Book Program,
  Reading, Mass., 1978.
\newblock {\sc doi: }\href{https://doi.org/10.1090/chel/364}{10.1090/chel/364}.

\bibitem{Al09}
L.~J.~S. Allen.
\newblock An introduction to stochastic epidemic models.
\newblock {\em Lecture Notes in Mathematics book series}, 6(8):1235--1252,
  2009.
\newblock {\sc doi:
  }\href{https://doi.org/10.1007/978-3-540-78911-6_3}{10.1007/978-3-540-78911-6\_3}.

\bibitem{Ar03}
M.~Arat\'o.
\newblock {A famous nonlinear stochastic equations (Lotka--Volterra model with
  diffusion)}.
\newblock {\em Math. Comp. Model.}, {\bf 38}:709--726, 2003.
\newblock {\sc doi:
  }\href{https://doi.org/10.1016/S0895-7177(03)90056-2}{10.1016/S0895-7177(03)90056-2}.

\bibitem{Ar74}
L.~Arnold.
\newblock {\em {Stochastic Differential Equations}}.
\newblock John Wiley \& Sons, New York, 1974.

\bibitem{Bz_14}
C.~Bai and H.~Zhang.
\newblock {A weak energy–momentum method for stochastic instability induced
  by dissipation and random excitations}.
\newblock {\em Probabilistic Eng. Mech.}, {\bf 37}:35--40, 2014.
\newblock {\sc doi:
  }\href{https://doi.org/10.1016/j.probengmech.2014.03.009}{10.1016/j.probengmech.2014.03.009}.

\bibitem{BBHLS15}
A.~Ballesteros, A.~Blasco, F.~J. Herranz, J.~de~Lucas, and C.~Sard\'{o}n.
\newblock Lie-{H}amilton systems on the plane: properties, classification and
  applications.
\newblock {\em J. Differential Equations}, 258(8):2873--2907, 2015.
\newblock {\sc doi:
  }\href{https://doi.org/10.1016/j.jde.2014.12.031}{10.1016/j.jde.2014.12.031}.

\bibitem{BM10}
D.~Bl\'{a}zquez-Sanz and J.~J. Morales-Ruiz.
\newblock Differential {G}alois theory of algebraic {L}ie-{V}essiot systems.
\newblock In {\em Differential algebra, complex analysis and orthogonal
  polynomials}, volume 509 of {\em Contemp. Math.}, pages 1--58. Amer. Math.
  Soc., Providence, RI, 2010.
\newblock {\sc doi:
  }\href{https://doi.org/10.1090/conm/509/09973}{10.1090/conm/509/09973}.

\bibitem{BM10local-global}
D.~Bl\'{a}zquez-Sanz and J.~J. Morales-Ruiz.
\newblock {Local and global aspects of Lie superposition theorem}.
\newblock {\em J. Lie Theory}, 20(3):483--517, 2010.
\newblock {\sc arXiv: }\href{https://arxiv.org/abs/0901.4478}{0901.4478}.

\bibitem{BCC_01}
J.~F. Blowey, J.~P. Coleman, and A.~W. Craig, editors.
\newblock {\em Theory and numerics of differential equations}.
\newblock Universitext. Springer-Verlag, Berlin, 2001.
\newblock Papers from the 9th EPSRC Numerical Analysis Summer School held at
  the University of Durham, Durham, July 10--21, 2000.

\bibitem{CCM19}
S.~Cai, Y.~Cai, and X.~Mao.
\newblock A stochastic differential equation {SIS} epidemic model with two
  independent {B}rownian motions.
\newblock {\em J. Math. Anal. Appl.}, 474(2):1536--1550, 2019.

\bibitem{CCM21}
S.~Cai, Y.~Cai, and X.~Mao.
\newblock A stochastic differential equation {SIS} epidemic model with regime
  switching.
\newblock {\em Discrete Contin. Dyn. Syst. Ser. B}, 26(9):4887--4905, 2021.

\bibitem{CFH23}
R.~Campoamor-Stursberg, E.~Fernández-Saiz, and F.~Herranz.
\newblock {Exact solutions and superposition rules for Hamiltonian systems
  generalising stochastic SIS epidemic models with variable infection rates}.
\newblock {\sc arXiv: }\href{https://arxiv.org/abs/2304.08843}{2304.08843}.

\bibitem{CL09}
J.~F. Cari\~{n}ena and J.~de~Lucas.
\newblock Quantum {L}ie systems and integrability conditions.
\newblock {\em Int. J. Geom. Methods Mod. Phys.}, 6(8):1235--1252, 2009.
\newblock {\sc doi:
  }\href{https://doi.org/10.1142/S021988780900420X}{10.1142/S021988780900420X}.

\bibitem{CL15}
J.~F. Cari\~{n}ena and J.~de~Lucas.
\newblock Quasi-{L}ie families, schemes, invariants and their applications to
  {A}bel equations.
\newblock {\em J. Math. Anal. Appl.}, 430(2):648--671, 2015.
\newblock {\sc doi:
  }\href{https://doi.org/10.1016/j.jmaa.2015.04.080}{10.1016/j.jmaa.2015.04.080}.

\bibitem{CLR11}
J.~F. Cari\~{n}ena, J.~de~Lucas, and M.~F. Ra\~{n}ada.
\newblock A geometric approach to integrability of {A}bel differential
  equations.
\newblock {\em Internat. J. Theoret. Phys.}, 50(7):2114--2124, 2011.
\newblock {\sc doi:
  }\href{https://doi.org/10.1007/s10773-010-0624-7}{10.1007/s10773-010-0624-7}.

\bibitem{CGL19}
J.~F. Cari\~{n}ena, J.~Grabowski, and J.~de~Lucas.
\newblock Quasi-{L}ie schemes for {PDE}s.
\newblock {\em Int. J. Geom. Methods Mod. Phys.}, 16(7):1950096, 36, 2019.
\newblock {\sc doi:
  }\href{https://doi.org/10.1142/S0219887819500968}{10.1142/S0219887819500968}.

\bibitem{CGM00}
J.~Cari{\~n}ena, J.~Grabowski, and G.~Marmo.
\newblock {\em {Lie--Scheffers Systems: A Geometric Approach}}, volume~3.
\newblock Napoli Series on Physics and Astrophysics, 2000.

\bibitem{CGM07}
J.~Cari{\~n}ena, J.~Grabowski, and G.~Marmo.
\newblock {Superposition rules, Lie theorem, and partial differential
  equations}.
\newblock {\em Rep. Math. Phys.}, {\bf 60}(2):237--258, 2007.
\newblock {\sc doi:
  }\href{https://doi.org/10.1016/S0034-4877(07)80137-6}{10.1016/S0034-4877(07)80137-6}.

\bibitem{Dissertationes}
J.~Cariñena and J.~de~Lucas.
\newblock {Lie systems: theory, generalisations, and applications}.
\newblock {\em Dissertationes Math.}, {\bf 479}:1--162, 2011.
\newblock {\sc doi:
  }\href{https://doi.org/10.4064/dm479-0-1}{10.4064/dm479-0-1}.

\bibitem{CLW22}
J.~Cariñena, J.~de~Lucas, and D.~Wysocki.
\newblock {Stratified Lie systems: theory and applications}.
\newblock {\em J. Phys. A: Math. Theor.}, {\bf 55}(38):385206, 2022.
\newblock {\sc doi:
  }\href{https://doi.org/10.1088/1751-8121/ac89bd}{10.1088/1751-8121/ac89bd}.

\bibitem{LS21}
J.~de~Lucas and C.~Sardón.
\newblock {\em A guide to Lie systems with compatible geometric structures}.
\newblock World Scientific, Singapore, 2020.
\newblock {\sc doi: }\href{https://doi.org/10.1142/q0208}{10.1142/q0208}.

\bibitem{LV15}
J.~de~Lucas and S.~Vilari\~{n}o.
\newblock {$k$}-symplectic {L}ie systems: theory and applications.
\newblock {\em J. Differential Equations}, 258(6):2221--2255, 2015.

\bibitem{LZ21}
J.~de~Lucas and B.~M. Zawora.
\newblock A time-dependent energy-momentum method.
\newblock {\em J. Geom. Phys.}, 170:Paper No. 104364, 22, 2021.

\bibitem{E82}
K.~Elworthy.
\newblock {\em {Stochastic differential equations on manifolds}}, volume~70 of
  {\em London Mathematical Society. Lecture notes series}.
\newblock Cambridge Univ. Press, 1982.

\bibitem{EFSZ22}
O.~Esen, E.~Fernández-Saiz, C.~Sardón, and M.~{Zaj\k ac}.
\newblock {A generalization of a SIS epidemic model with fluctuations}.
\newblock {\em Math. Methods Appl. Sci.}, {\bf 45}:3718--3731, 2022.
\newblock {\sc doi: }\href{https://doi.org/10.1002/mma.8013}{10.1002/mma.8013}.

\bibitem{Ev13}
L.~C. Evans.
\newblock {\em An introduction to stochastic differential equations}.
\newblock American Mathematical Society, Providence, RI, 2013.

\bibitem{GGHMP11}
A.~Gray, D.~Greenhalgh, L.~Hu, X.~Mao, and J.~Pan.
\newblock {A stochastic differential equation SIS epidemic}.
\newblock {\em SIAM J. Appl. Math.}, {\bf 71}(3):876--902, 2011.
\newblock {\sc doi:
  }\href{https://doi.org/10.1137/10081856X}{10.1137/10081856X}.

\bibitem{Gu93}
A.~Guldberg.
\newblock {Sur les équations différentielles ordinaires qui possèdent un
  système fondamental d’intégrales}.
\newblock {\em C.R. Math. Acad. Sci. Paris}, {\bf 116}:964--965, 1893.

\bibitem{HH11}
R.~Habibi, H.~Habibi, and H.~Habibi.
\newblock Three notes about linear stochastic oscillator equations.
\newblock {\em Int. Math. Forum}, 6(61-64):3131--3137, 2011.

\bibitem{Ha19}
S.~Hafstein.
\newblock Lyapunov functions for linear stochastic differential equations:
  {BMI} formulation of the conditions.
\newblock In {\em I{CINCO} 2019---{P}roceedings of the 16th {I}nternational
  {C}onference on {I}nformatics in {C}ontrol, {A}utomation and {R}obotics.
  {V}ol. {I}}, pages 147--155. SciTePress, Set\'{u}bal, [2019] \copyright 2019.
\newblock {\sc doi:
  }\href{https://doi.org/10.5220/0008192201470155}{10.5220/0008192201470155}.

\bibitem{LO09}
J.~L\'azaro-Cam\'{i} and J.~Ortega.
\newblock {Superposition rules and stochastic Lie--Scheffers systems}.
\newblock {\em Ann. Inst. Henri Poincaré Probab. Stat.}, {\bf 45}(4):910--931,
  2009.
\newblock {\sc doi:
  }\href{https://doi.org/10.1214/08-AIHP189}{10.1214/08-AIHP189}.

\bibitem{LO08}
J.-A. L\'{a}zaro-Cam\'{\i} and J.-P. Ortega.
\newblock Stochastic {H}amiltonian dynamical systems.
\newblock {\em Rep. Math. Phys.}, 61(1):65--122, 2008.
\newblock {\sc doi:
  }\href{https://doi.org/10.1016/S0034-4877(08)80003-1}{10.1016/S0034-4877(08)80003-1}.

\bibitem{LP08}
J.-A. L\'{a}zaro-Cam\'{\i} and J.-P. Ortega.
\newblock Stochastic {H}amiltonian dynamical systems.
\newblock {\em Rep. Math. Phys.}, 61(1):65--122, 2008.
\newblock {\sc doi:
  }\href{https://doi.org/10.1016/S0034-4877(08)80003-1}{10.1016/S0034-4877(08)80003-1}.

\bibitem{LO9Red}
J.-A. L\'{a}zaro-Cam\'{\i} and J.-P. Ortega.
\newblock Reduction, reconstruction, and skew-product decomposition of
  symmetric stochastic differential equations.
\newblock {\em Stoch. Dyn.}, 9(1):1--46, 2009.
\newblock {\sc doi:
  }\href{https://doi.org/10.1142/S0219493709002531}{10.1142/S0219493709002531}.

\bibitem{Li93a}
M.~Lie.
\newblock {On differential equations possessing fundamental integrals}.
\newblock {\em Leipziger Berichte}, 1893.

\bibitem{Lie1893}
S.~Lie.
\newblock {\em {Vorlesungen über continuierliche Gruppen mit geometrischen und
  anderen Anwendungen}}.
\newblock Teubner, 1893.
\newblock \href{http://eudml.org/doc/203682}{http://eudml.org/doc/203682}.

\bibitem{Ma78}
P.~Malliavin.
\newblock {\em G\'{e}om\'{e}trie diff\'{e}rentielle stochastique}, volume~64 of
  {\em S\'{e}minaire de Math\'{e}matiques Sup\'{e}rieures [Seminar on Higher
  Mathematics]}.
\newblock Presses de l'Universit\'{e} de Montr\'{e}al, Montreal, Que., 1978.
\newblock Notes prepared by Dani\`ele Dehen and Dominique Michel.

\bibitem{MS88}
J.~E. Marsden and J.~C. Simo.
\newblock The energy momentum method.
\newblock {\em Act. Acad. Sci. Tau.}, 1(124):245--268, 1988.

\bibitem{NM19}
G.~Nakamura and A.~Mart\'inez.
\newblock {Hamiltonian dynamics of the SIS epidemic model with stochastic
  fluctuations}.
\newblock {\em Scientific Reports}, {\bf 9}:15841, 2019.
\newblock {\sc doi:
  }\href{https://doi.org/10.1038/s41598-019-52351-x}{10.1038/s41598-019-52351-x}.

\bibitem{Ot21}
O.~M. Otunuga.
\newblock Time-dependent probability distribution for number of infection in a
  stochastic {SIS} model: case study {COVID}-19.
\newblock {\em Chaos Solitons Fractals}, 147:Paper No. 110983, 20, 2021.

\bibitem{Pi12}
G.~Pietrzkowski.
\newblock {Explicit solutions of the {${\mathfrak{a}_1}$}-type Lie--Scheffers
  system and a general Riccati equation}.
\newblock {\em J. Dyn. Control Syst.}, 18(4):551--571, 2012.
\newblock {\sc doi:
  }\href{https://doi.org/10.1007/s10883-012-9159-y}{10.1007/s10883-012-9159-y}.

\bibitem{Sa70}
P.~Sagirow.
\newblock {\em {Stochastic Methods in the Dynamics of Satellites }}.
\newblock CISM International Centre for Mechanical Sciences. UDINE, 1970.

\bibitem{Ka81}
N.~G. van Kampen.
\newblock It\^{o} versus {S}tratonovich.
\newblock {\em J. Statist. Phys.}, 24(1):175--187, 1981.

\bibitem{Ve93}
M.~Vessiot.
\newblock {Sur une classe d'équations différentielles}.
\newblock {\em Ann. Sci. École Norm. Sup.}, {\bf 10}:53--64, 1893.

\bibitem{WCDG18}
W.~Wang, Y.~Cai, Z.~Ding, and Z.~Gui.
\newblock A stochastic differential equation {SIS} epidemic model incorporating
  {O}rnstein-{U}hlenbeck process.
\newblock {\em Phys. A}, 509:921--936, 2018.

\bibitem{PW83}
P.~Winternitz.
\newblock Lie groups and solutions of nonlinear differential equations.
\newblock In {\em Nonlinear phenomena ({O}axtepec, 1982)}, volume 189 of {\em
  Lecture Notes in Phys.}, pages 263--331. Springer, Berlin, 1983.
\newblock {\sc doi:
  }\href{https://doi.org/10.1007/3-540-12730-5_12}{10.1007/3-540-12730-5\_12}.

\end{thebibliography}
